\documentclass[sigconf, review]{acmart}

\newcommand{\bbd}[1]{\textbf{\textcolor{blue}{#1}}}
\newcommand{\obd}[1]{\textbf{\textcolor{orange}{[TODO: #1]}}}
\newcommand{\rbd}[1]{\textbf{\textcolor{red}{[QUESTION: #1]}}}
\newcommand{\com}[1]{{\small \textcolor{gray}{[COMMENT: #1]}}}

\newcommand{\revone}[1]{\textcolor{red}{#1}}
\newcommand{\revtwo}[1]{\textcolor{orange}{#1}}
\newcommand{\revthr}[1]{\textcolor{blue}{#1}}

\usepackage{lipsum}
\usepackage{enumitem}
\usepackage{changepage}
\usepackage[T1]{fontenc}
\usepackage{hyperref}
% \usepackage[margin=1in]{geometry}
\usepackage{graphicx} 
\usepackage{xcolor}
% \usepackage{amsmath}
% \usepackage{amsthm}
% \usepackage{amssymb}
% \usepackage{newtxtext,newtxmath}
% \usepackage{algorithm}
\usepackage[ruled,linesnumbered]{algorithm2e}
\DontPrintSemicolon
% \usepackage{algpseudocode} 

% \usepackage{algorithm2e}

% \usepackage{setspace} 
\usepackage{tikz}
\usetikzlibrary{mindmap, trees}
% \onehalfspacing 
% \algrenewcommand\algorithmiccomment[1]{\hfill\textit{#1}}
\SetKw{KwAnd}{and}


\SetKwInput{KwParam}{Public Parameters}
\theoremstyle{definition} 
\theoremstyle{definition}  % 设置定义的样式
\newtheorem{theorem}{Theorem}[section] 
\newtheorem{definition}{Definition}[section]
\newtheorem{proposition}{Proposition}[section]
\newtheorem{lemma}{Lemma}[section]
\newtheorem{corollary}{Corollary}[section] 
\newtheorem*{observation}{Observation}
\newtheorem{example}{Example}[section]
\newtheorem*{remark}{Remark}



\renewcommand{\today}{\number\day\ \ifcase\month\or January\or February\or March\or April\or May\or June\or July\or August\or September\or October\or November\or December\fi\ \number\year}

\newcommand{\err}{\mathrm{Error}} 



\AtBeginDocument{
  \providecommand\BibTeX{{
    Bib\TeX}}}


\setcopyright{acmlicensed}
\copyrightyear{2018}
\acmYear{2018}
\acmDOI{XXXXXXX.XXXXXXX}

\acmConference[CCS '26]{The 2026 ACM SIGSAC Conference on Computer and Communications Security}
  {November 15--19, 2026}{The Hague, The Netherlands}

\acmISBN{978-1-4503-XXXX-X/2018/06}



\begin{document}

\title{Domain Partitioning for Shuffle-DP}
\author{Anonymous Author(s)}




\begin{abstract}
    % \lipsum[0-2]
\end{abstract}


\begin{CCSXML}
<ccs2012>
   <concept>
       <concept_id>10002978</concept_id>
       <concept_desc>Security and privacy</concept_desc>
       <concept_significance>500</concept_significance>
       </concept>
 </ccs2012>
\end{CCSXML}

\ccsdesc[500]{Security and privacy}

%%
%% Keywords. The author(s) should pick words that accurately describe
%% the work being presented. Separate the keywords with commas.
\keywords{Differential privacy}
%% A "teaser" image appears between the author and affiliation
%% information and the body of the document, and typically spans the
%% page.


%%
%% This command processes the author and affiliation and title
%% information and builds the first part of the formatted document.
\maketitle





\section{Introduction}

% 传统DP研究了很多年，有两种DP setting：record和user。user-DP更贴合现实场景，因为一个user可能....。

% highlight一下user-DP重要性

% 再motivate一下shuffle-DP， highlight一下。参考上一片的写法，在local和central之间 trade-off

% 然后说一下shuffle-DP已经在很多问题比如counting，graph，xxx很好的研究。传统的shuffle-DP集中在static setting，但我们注意到已经有研究dynamic setting了，但是都还是集中在record-DP上。给个例子说user-DP很常见（应该是给个场景）。

% 我们是第一篇研究user-shuffle-DP的。gap


\section{Preliminary}
\label{sec:Preliminary}


This section establishes the technical foundations of our methodology.
We begin by formalizing the notation in Section~\ref{sec:Notation}, review fundamental DP concepts and mechanisms in Section~\ref{sec:DP Basics}, and formally define the shuffle-DP model in Section~\ref{sec:Shuffle-DP}.



\subsection{Notation}
\label{sec:Notation}


Let $[U] = \{0, 1, \dots, U\}$ and $\mathcal{X} = [U]^d$ denote the input domain. 
Let $n$ denote the number of users. 
The dataset contributed by user $i \in [n]$ is $D_i = (x_{i,1}, x_{i,2}, \ldots, x_{i,m_i})$, where $m_i \leq M$ denotes the number of records contributed by user $i$ and $M$ is the maximum number of records allowed per user. 
The entire dataset is $D = \biguplus_{i=1}^n D_i$.
Let $| \cdot |$ denote the $\ell_1$-norm for any vector; then $| D_i | = m_i$.
The entire dataset is then $D = \bigcup_{i=1}^n D_i$.
Given a query $Q$, the true query result is $Q(D)$ and its estimate is $\tilde{Q}(D)$. 
We define $d(D, D')$ as the distance between datasets $D$ and $D'$, measured by the minimum number of single-record changes required to transform $D$ into $D'$.


% \begin{definition}[Record-Level Neighbors]
%     Two datasets $D$ and $D'$ are record-level neighbors, denoted as $D \sim_r D'$, if $d(D, D') = 1$, i.e., $D' = D \cup \{x\} $ or $ D' = D \setminus \{x\}. $ 
% \end{definition}


\begin{definition}[Record-Level Neighbors]
Two datasets $D$ and $D'$ are record-level neighbors, denoted $D \sim_r D'$, if $d(D, D') = 1$, i.e., there exists $x \in D$ and $x'\in D'$ such that $x \neq x'$ while all other records are identical.
\end{definition}


% please check: no need of $d()$?

% \begin{definition}[User-Level Neighbors]
%     Two datasets $D$ and $D'$ are user-level neighbors, denoted as $D \sim_u D'$, if they differ only by the entire contribution $D_i$ of exactly one user $i$, i.e., $D' = D \cup D_i $ or $ D' = D \setminus D_i.$
% \end{definition}

\begin{definition}[User-Level Neighbors]
Two datasets $D$ and $D'$ are user-level neighbors, denoted $D \sim_u D'$, if one can be obtained from the other by replacing all or any subset of records contributed by exactly one user $i$, while all other users' records remain identical.
\end{definition}



In the dynamic setting, time proceeds in discrete steps, and data is modeled as a stream of items arriving over time, one per time step. We focus on the finite time horizon setting where the stream length is $T$, i.e., $t \in [T]$.
For each user $i$, let $D_i^t = (x_{i,1}, x_{i,2}, \ldots, x_{i,t})$ denote user $i$'s data stream up to time $t$, where $x_{i,s} \in \mathcal{X} \cup \{\perp\}$ and $x_{i,s} = \perp$ indicates no contribution at time $s$. 
The entire dataset at time $t$ is $D^t = \biguplus_{i=1}^n D_i^t$, and $Q(D^t)$ denotes the true query result at time $t$. 
A mechanism for the dynamic setting must output an estimate $\tilde{Q}(D^t)$ at each time step $t$, before the arrival of the next record at time $t+1$.




% please check, whether we need to define "or simple event-DP"

Following~\cite{dwork2010differential}, we distinguish two neighboring relations in the dynamic setting. 
In \textit{event-level DP}, each user contributes at most one real record throughout the stream, i.e., $|\{t \in [T] : x_{i,t} \neq \perp\}| = 1$, and two datasets are neighbors if they differ in a single event: there exists user $i$ and time $t$ such that $x_{i,t} \neq x'_{i,t}$ while $x_{j,s} = x'_{j,s}$ for all $(j,s) \neq (i,t)$. 
In \textit{user-level DP}, each user may contribute up to $M$ records, and two datasets are neighbors if they differ in all or any subset of events of exactly one user: there exists user $i$ such that $D_i^T \neq D_i^{\prime T}$ while $D_j^T = D_j^{\prime T}$ for all $j \neq i$.




% In the dynamic setting, time proceeds in discrete steps. At each time step $t \in \{1,2,\dots,T\}$, a new record $x_t\in \{\perp\} \cup \mathcal{X}$ is observed, where \(x_t = \perp\) represents that no user contributes a record at time \(t\). The data stream is thus represented as $D=(x_1, x_2, \dots, x_T).$
% At time step $t$, the entire dataset containing all records observed so far is $ D^t = (x_1, x_2, \dots, x_t)$. 
% For each user $i$, we denote by $D_i^t = (x_{i,1}, x_{i,2}, \ldots, x_{i,m_i^t})$ the subset of records contributed by user $i$ up to time $t$, where $m_i^t$ is the cumulative number of records from user $i$ observed so far. 
% Consequently, the entire dataset at time $t$ can also be expressed as $D^t = \bigcup_{i=1}^n D_i^t$.
% We denote by $Q(D^t)$ the true query result at time $t$. The mechanism then outputs an estimate $\tilde{Q}(D^t)$, computed prior to the arrival of the next record at time $ t+1 $.
% Following~\cite{dwork2010differential}, we distinguish two types of neighboring relationships in the dynamic setting:
% in \textit{event-level DP}, two data streams $D$ and $D'$ are neighbors if they differ by a single event (i.e., $d(D, D')=1$),
% whereas in \textit{user-level DP}, two streams are neighbors if they differ in all events contributed by a single user.

% dynamic neighbors?


% The query function $Q$ in our framework can be both union-preserving and non–union-preserving queries. Some common queries are listed below:
We illustrate several common queries $Q$ defined on a dataset $D$:

\begin{enumerate}
    % \item Bit counting: $Q_{\text{count}}(D) = \sum_{x \in D} x$;
    \item Summation: $Q_{\text{sum}}(D) = \sum_{x \in D} x$;
    \item Frequency estimation: $Q_{\text{hist}}(D) = ( \sum_{x \in D} \mathbb{I}[x = j] )_{j=0}^{U}$; % where $\mathcal{X} = \{0, 1, \cdots, U\}$;
    \item Range counting: $Q_{\text{range}}(D) = ( \sum_{x \in D}\mathbb{I}[\,l \le x \le r\,] )_{0 \leq l \leq r < U}$; 
    \item Median: 
    \[
    \begin{aligned}
        Q_{\mathrm{median}}(D) = \frac{1}{2} (
        & \min\{ v\in D \Bigm| \sum_{x \in D}\mathbb{I}[x\le v]\ge\frac{n}{2} \} \\
        + \  &\max\{ v\in D \Bigm| \sum_{x \in D}\mathbb{I}[x< v]\le\frac{n}{2}\}).
    \end{aligned}
    \]
\end{enumerate}

Note that the bit counting query $Q_\text{count}$ can be viewed as a specific instance of the summation query $Q_\text{sum}$ with $U=1$. 
Let $\mathrm{\gamma}_{\ell_p}(Q,m)$ denote the maximum change in the query result (in $\ell_p$-norm) when at most $m$ records are removed:
\[\mathrm{\gamma}_{\ell_p}(Q,m) =\max \big\{\|Q(D)-Q(D')\|_{\ell_p}:\ d(D,D')=m\big\}.\]
For example, $\gamma_{\ell_1}(Q_\text{sum},m)=m\cdot U.$


\subsection{DP Basics}
\label{sec:DP Basics}
Below we review the definition of DP, several standard DP properties, and two fundamental DP mechanisms.
Unless otherwise specified, we follow the standard record-level DP setting to introduce these concepts.


\begin{definition}[Differential Privacy]
\label{def: DP}
    For some $\varepsilon > 0$ and $0 \leq \delta < n^{-\Omega(1)}$, a randomized mechanism $\mathcal{M}:\mathcal{X}^{n} \to \mathcal{Y}$ is $(\varepsilon, \delta)$-differentially private if for any two neighboring datasets $D \sim_r D'$, $\mathcal{M}(D)$ and $\mathcal{M}(D')$ are $(\varepsilon, \delta)$-indistinguishable. That is, for any subset of outputs $Y\subseteq \mathcal{Y}$,
    \[
    \Pr[\mathcal{M}(D) \in Y ] \leq e^\varepsilon \cdot \Pr[\mathcal{M}(D') \in Y ] + \delta.
    \]
\end{definition}



In practice, $\varepsilon$ is usually a constant between $0.1$ and $10$, and $\delta$ should be much smaller than $1/n$.
Moved to the dynamic setting, a randomized mechanism $\mathcal{M}$ is $(\varepsilon, \delta)$-differentially private if for any two neighboring datasets $D$ and $D'$, the entire output is $(\varepsilon, \delta)$-indistinguishable. That is, for any $Y_1, Y_2, \dots, Y_T \subseteq \mathcal{Y}$,
\[
\begin{aligned}
&\Pr\!\big[ \mathcal{M}(D^1)\in Y_1, \mathcal{M}(D^2)\in Y_2, \dots, \mathcal{M}(D^T)\in Y_T \big] \\
&\le\;
e^{\varepsilon}\cdot
\Pr\!\big[ \mathcal{M}(D'^1)\in Y_1,\ \mathcal{M}(D'^2)\in Y_2, \dots, \mathcal{M}(D'^T)\in Y_T \big]
+\delta .
\end{aligned}
\]


Below are some commonly used DP properties.

\begin{lemma} [Basic Composition~\cite{dwork2014algorithmic}]
\label{lem:Sequential Composition}
If $ \mathcal{M} $ is a (possibly adaptive) composition of differentially private mechanisms $ \mathcal{M}_1,\mathcal{M}_2,\dots, \mathcal{M}_k$, where each $ \mathcal{M}_i$ satisfies $(\varepsilon, \delta)$-DP, then $ \mathcal{M} $ satisfies $(k\varepsilon, k\delta)$-DP.
\end{lemma} 

\begin{lemma} [Parallel Composition~\cite{mcsherry2009privacy}]
\label{lem:Parallel Composition}
    Let $ \mathcal{X}_1, \dots, \mathcal{X}_k $ be pairwise disjoint subdomains of $ \mathcal{X} $, and each $ \mathcal{M}_i : \mathcal{X}^n_i \to \mathcal{Y} $ be an $(\varepsilon, \delta)$-DP mechanism. Then the mechanism $ \mathcal{M}(D) := (\mathcal{M}_1\left(D \cap \mathcal{X}_1 \right)$, $\dots$, $\mathcal{M}_k \left(D \cap \mathcal{X}_k \right))$ satisfies $(\varepsilon, \delta)$-DP.
\end{lemma}

\begin{lemma} [Group Privacy~\cite{dwork2014algorithmic}]
\label{lem:Group Privacy}
    If $\mathcal{M}$ is an $\varepsilon$-DP mechanism, then for any two datasets $D$, $D'$ with $d(D,D')=k$, $\mathcal{M}$ satisfies $(k\varepsilon)$-DP.
\end{lemma}

\subsubsection{Laplace mechanism} 
A basic DP mechanism is the Laplace mechanism. Given any query $Q$, the \textit{global sensitivity} is defined as $\mathrm{GS}_Q = \max_{D\sim_rD'}||Q(D)-D(D')||$.
\begin{lemma}[Laplace Mechanism]
\label{lem:Laplace Mechanism}
    Given any query $Q$, the mechanism $\mathcal{M}(D) = Q(D) + \eta$ preserves $\varepsilon$-DP, where $\eta$ is drawn from a Laplace distribution with scale $\mathrm{GS}_Q/\varepsilon$, i.e., $\eta \sim \mathrm{Lap}(\mathrm{GS}_Q/\varepsilon)$. With high probability at least $1-\beta$, the error of the Laplace mechanism is at most $\textrm{GS}_Q/\varepsilon\cdot \log (1/\beta)$.
\end{lemma}
% Please check if there can be a log? may should be ln

\begin{lemma}[Concentration Bound of Laplace Distributions~\cite{chan2011private}]
\label{lem:Concentration Bound}
    Suppose $\eta_1,\eta_2,\dots,\eta_k$ are independent random variable, where each $\eta_i \sim \mathrm{Lap}(b_i)$, then for any $\beta>0$:
    \[
    \Pr\left[ \left| \sum_i\eta_i \right| \le \sqrt{8\sum_i b_i^2} \cdot \log (\frac{2}{\beta})\right] \le \beta.
    \]
\end{lemma}

As shown in Lemma~\ref{lem:Laplace Mechanism}, in the static setting, the Laplace mechanism returns an estimate $\tilde{Q}(D)$ with error $O(\textrm{GS}_Q/\varepsilon)$. 
% Please check which word we should use to describe the data/record/event in the stream?
However, when we move to the dynamic setting and apply the Laplace mechanism independently to each $x_t$ at time $t$, the estimate $\tilde{Q}(D^t) = \sum_{t=1}^T \tilde{Q}(x_t)$ incurs a total error that grows as $O(\sqrt{t} \cdot \textrm{GS}_Q/\varepsilon)$ by the concentration bound (Lemma~\ref{lem:Concentration Bound}), resulting in poor utility.
To overcome this limitation, we introduce the binary mechanism which achieves a polylogarithmic error in $t$.

\subsubsection{Binary mechanism}
\label{sec:Binary mechanism}
We briefly introduce the binary mechanism (or continual observation mechanism) under event-DP~\cite{dwork2010differential, chan2011private}, which will be used in our framework. 
The core idea is to apply a binary decomposition over the $T$ time steps, constructing $\log T$ levels of intervals, and the Laplace mechanism is applied to each interval to produce a noisy count.
To compute $\tilde{Q}_\text{count}(D^t)$ at time $t$, the mechanism selects at most $\log T$ intervals that form a non-overlapping cover of $[1,t]$. 
Summing the noisy counts of these intervals yields the estimate at time $t$.
To ensure overall DP, each interval is allocated a privacy budget of $\varepsilon / \log T$, which is justified by basic composition (Lemma~\ref{lem:Sequential Composition}) across levels and parallel composition (Lemma~\ref{lem:Parallel Composition}) within levels.

\begin{lemma}
\label{lem:binary mechanism}
    Given any $\varepsilon > 0 $, for any $D$, any $t \in \{1,2,\dots,T\}$, and any $\beta> 0$, with probability at least $1 - \beta$, the binary mechanism returns an $\tilde{Q}_\text{count}(D^t)$ such that
    \[
    \left| \tilde{Q}_\text{count}(D^t)-{Q}_\text{count}(D^t) \right| \le \frac{4}{\varepsilon} \cdot \lceil \log(t) \rceil ^ {1.5} \cdot \log \frac{1}{\beta}.
    \]
\end{lemma}


% in the finite stream, replace beta with beta/2T (t^2 because it's infinite stream)
\begin{remark}
    Based on the one-shot guarantee (Lemma~\ref{lem:binary mechanism}), an all-time guarantee follows by replacing the failure probability $\beta$ with $\beta/T$. The total failure probability is at most $\sum_{t=1}^T \beta/T = O(\beta)$ by union bound, and the $\log(1/\beta)$ term in the error bound becomes $O(\log(T/\beta))$.
\end{remark}

% Compared with the $O(t)$ error obtained by independently applying the Laplace mechanism at each time step, the binary mechanism achieves a polylogarithmic error in $t$.

% \begin{remark}
%     Although originally proposed for central-DP, the binary mechanism extends naturally to the shuffle-DP setting as it fundamentally relies on a combination of multiple counting queries.
% \end{remark}

Although originally proposed for central-DP, the binary mechanism extends naturally to the shuffle-DP setting. 
The key observation is that each node in the binary tree corresponds to a counting query over a specific time interval, which can be implemented using shuffle-DP protocols.
In this paper, we replace the Laplace mechanism at each tree node with~\cite{ghazi2021differentially}, which will be introduced later in Example~\ref{exm:Bit counting problem}.
\begin{lemma}
\label{lem:shuffle binary mechanism}
Given any $\varepsilon > 0$ and $\beta > 0$, with probability at least $1-\beta$, for all time steps $t \in \{1,2,\dots,T\}$ simultaneously,  the binary mechanism under the shuffle-DP protocol~\cite{ghazi2021differentially} achieves an error of $O(\tfrac{\log^{1.5}T}{\varepsilon}\log\tfrac{T}{\beta})$ for the dynamic bit counting problem.
\end{lemma}

% pleach check here: do we need a lemma?



\subsection{Shuffle-DP}
\label{sec:Shuffle-DP}
Different DP models can be captured by Definition~\ref{def: DP} by the formulation of $\mathcal{M}(D)$.
In central-DP, $\mathcal{M}(D)$ represents the output of a trusted curator. In contrast, for local-DP and shuffle-DP, where the analyzer $\mathcal{A}$ is untrusted, $\mathcal{M}(D)$ is defined by the view of $\mathcal{A}$.
Under local-DP, each user $i$ independently privatizes its own data $x_i$ using a local randomizer $\mathcal{R}$ and then sends the randomized response $\mathcal{R}(x_i)$ to $\mathcal{A}$ hence $\mathcal{M}(D)=(\mathcal{R}(x_1), \mathcal{R}(x_2), \dots, \mathcal{R}(x_n))$.
Under shuffle-DP, a trusted shuffler $\mathcal{S}$ randomly permutes the responses before forwarding them to analyzer $\mathcal{A}$, resulting in
\[\mathcal{M}(D) = \mathcal{S} \circ \mathcal{R} (D)= \{\mathcal{R}(x_1), \mathcal{R}(x_2), \dots, \mathcal{R}(x_n)\}.\]

Under shuffle-DP, a protocol which is used to answer query $Q$ is defined as $\mathcal{P}_Q = (\mathcal{R}, \mathcal{S}, \mathcal{A})$, where $\mathcal{R}, \mathcal{S}, \mathcal{A}$ denote the randomizer, shuffler, and analyzer, respectively. The protocol output is given by
\[\mathcal{P}_Q(D)\gets \mathcal{A} \circ \mathcal{S} \circ \mathcal{R} (D) = \mathcal{A}(\mathcal{S}(\mathcal{R}(x_1), \cdots, \mathcal{R}(x_n))).\]
%$\varepsilon$ and $\delta$ are privacy budget and 
$\mathcal{S} \circ \mathcal{R} (D)$ satisfies $(\varepsilon,\delta)$-DP. 


\paragraph{Multiple shuffler setting.} 
Some shuffle-DP protocols consider the setting where there exist multiple shufflers $\mathcal{S}$~\cite{balle2020privatesum}.
In this paper, we have a set of shufflers $\{\mathcal{S}_1, \mathcal{S}_2, \ldots, \mathcal{S}_m\}$ and $n$ users. 
Each user $i$ can send messages to one or more shufflers $\mathcal{S}$. 


\paragraph{Shuffle-DP protocol parameters}
 By convention, a shuffle-DP protocol has four parameters: the privacy budgets $\varepsilon$ and $\delta$, the error parameter $\beta$, and the data size $n$. The parameter $\beta$ specifically controls the probability of exceeding the stated error bound. 
More precisely, we use $\err_{\ell_p}(\mathcal{P}_Q, \varepsilon, \delta, \beta,n)$ to denote the $\ell_p$-norm error of the protocol $\mathcal{P}_Q$. Specifically, for any input dataset $D$, with privacy budgets $\varepsilon, \delta$, with probability at least $1 - \beta$, we have:
\[
\left\| \mathcal{P}_Q(D) - Q(D) \right\|_{\ell_p} \leq \err_{\ell_p}(\mathcal{P}_Q, \varepsilon, \delta, \beta,n).
\]


We quantify the communication cost by $\mathrm{Msg}(\mathcal{P}_Q,\varepsilon,\delta,n)$, which denotes the expected number of messages each user sends in protocol $\mathcal{P}_Q$ with a group of $n$ users, i.e., $\mathbb{E}[|\mathcal{R}(x_i)|]$. 
For all existing shuffle-DP protocols, it is observed that $\err_{\ell_p}(\mathcal{P}_Q, \varepsilon, \delta, \beta,n)$ and $\mathrm{Msg}(\mathcal{P}_Q, \varepsilon, \delta, n)$ are both polynomial functions of their respective parameters. 
This observation will simplify the analysis in this paper. For example, $\err{\ell_p}(\mathcal{P}Q, \varepsilon/c, \delta/c, \beta/c,n) = O\big(\err{\ell_p}(\mathcal{P}_Q, \varepsilon, \delta, \beta,n)\big)$ for any constant $c$.




\begin{example}[Shuffle-DP Protocol $\mathcal{P}_{Q_{\text{count}}}$~\cite{ghazi2021differentially}] 
\label{exm:Bit counting problem} 
We take the bit counting problem ($Q_\text{count}$) under record-level DP as an example to illustrate how a shuffle-DP protocol works. 
Each user $i$ holds a bit \( x_i \in \{0, 1\} \), and we aim to compute the sum $\sum_ix_i$. 
The state-of-the-art shuffle-DP protocol~\cite{ghazi2021differentially} works as follows: each user sends a real $+1$'' message if $x_i = 1$. 
To preserve DP, each user independently samples a random number of dummy $+1$'' and ``$-1$'' messages from a Negative Binomial distribution. 
After shuffling, the noises contributed by all $n$ users aggregate to a distribution that approximates a discrete Laplace distribution, thereby achieving accuracy close to that of central-DP. 
The protocol of~\cite{ghazi2021differentially} has the following properties: 
\begin{lemma}
\label{lem: Bit Counting Protocol} 
Given $\varepsilon >0, \delta>0$ and $n$, the protocol $\mathcal{P}{Q{\text{count}}}$ of~\cite{ghazi2021differentially} solves the bit counting problem under $(\varepsilon,\delta)$-shuffle-DP with the following guarantees: 
\begin{itemize} 
\item Utility guarantee: $\err_{\ell_1}(\mathcal{P}{Q{\text{count}}}, \varepsilon, \delta, \beta,n)=O\left(\tfrac{1}{\varepsilon}\log\tfrac{1}{\beta}\right)$; 
\item Communication cost: $\mathrm{Msg}(\mathcal{P}_{Q_{\text{count}}}, \varepsilon, \delta, n)=1+{O}\left(\tfrac{\log(1/\delta)}{\varepsilon n}\right)$. 
\end{itemize} 
\end{lemma} 
\end{example}


\section{Naive Solutions and Objective}

% simple attempts, challenge 显得在凑页数

% new structure: Objectives and Challenges.
% first paragraph: core difficulty: shuffle-DP needs uniform contribution.
% challenge1 utility: use uniform M. -> low utility O(M) -> objective 1: we want a instance-specific error
% challenge2 privacy: use max_(D) to obtain instance-specific error -> break privacy -> objective 2: preserve DP
% in summary, ...

% please check that, the overall error may include the dummy message (you also introduce noise)

\label{sec:Challenges and Objectives}

In this section, we study the fundamental challenges of user-level shuffle-DP in the static setting. Before presenting our final solution, we analyze several naive approaches to highlight the challenges. These preliminary explorations provide valuable insights that directly guide the design of our framework.


The core challenge of user-level shuffle-DP lies in the natural variation in user contributions. 
The traditional shuffle-DP protocol requires each user to hold a single data element.
However, in the user-DP setting, the number of records $m_i$ contributed by each user $i$ can vary significantly. 
Although for simple additive queries such as $Q_{\text{sum}}$, we can effectively reduce the problem to record-DP setting by asking users to locally aggregate their records into a single record, this approach fails to generalize to queries such as $Q_{\text{hist}}$, $Q_{\text{range}}$, or $Q_{\text{median}}$, which require detailed distributional information that will be lost when compressing a user's dataset into a single record.
Generally, accommodating such heterogeneity poses fundamental challenges regarding both privacy and utility, which we discuss below.

\paragraph{Attempt 1: Global Padding}

To address the variation in $m_i$, a straightforward solution is to enforce a global upper bound $M$ on user contributions. 
Under this strategy, all users pad their datasets with dummy records to reach size $M$ and randomize each record using privacy budget $\varepsilon' = \varepsilon / M$. 
While this approach handles the variation in user contributions, it relies on the pre-defined bound $M$, which must be set conservatively to cover all possible inputs. 
In practice, user contributions may be far below $M$, leading to excessive noise and substantial utility degradation. 
Consider query $Q_\text{count}$ in a sparse scenario where every user contributes a single record (i.e., $m_i = 1$ for all $i$). 
The optimal noise scale is $O(1/\varepsilon)$, yet global padding forces every user to split their privacy budget by $M$, resulting in a noise scale $O(M/\varepsilon)$. 
A more ideal approach achieves instance-specific error depending on $m_{\max}(D)$, the maximum number of records contributed by any user in dataset $D$. 
In most inputs, $m_{\max}(D)$ can be substantially smaller than $M$ and better describes the actual data characteristics. 
Therefore, our first objective is to achieve instance-specific error depending on $m_{\max}(D)$ rather than a pre-defined global bound $M$.





\paragraph{Attempt 2: Clipping with $m_{\max}(D)$}



To achieve instance-specific error, we consider the \textit{clipping mechanism}~\cite{abadi2016deep,huang2021instance}, widely used in central-DP~\cite{dong2024almost} to reduce error dependency from a pre-defined global value to an instance-specific one. 
Applying this to user contributions, we select a threshold $m_\tau$ such that users contributing more than $m_\tau$ records truncate their datasets to size $m_\tau$, while users with fewer records pad their datasets with $\perp$ to reach exactly $m_\tau$. 
The protocol then runs a standard record-level shuffle-DP mechanism with privacy budget $\varepsilon/m_\tau$ on the standardized dataset, reducing noise dependency from $M$ to $m_\tau$. 
However, clipping introduces bias by discarding records from users where $m_i > m_\tau$. 
To minimize this bias, we seek $m_\tau \approx m_{\max}(D)$ so that only a small fraction of records are clipped. 
A naive approach is to set $m_\tau = m_{\max}(D)$ directly, but this violates DP since $m_{\max}(D)$ is itself a sensitive statistic revealing information about the most active user. 
Therefore, we must estimate a suitable clipping threshold $m_\tau$ in a DP manner.





\paragraph{Central-DP Solutions for Estimating $m_\tau$}

The problem of privately estimating a threshold $m_\tau \approx m_{\max}(D)$ is well-studied under central-DP, with two common approaches:


(1) \textit{Sparse Vector Technique (SVT)}~\cite{dwork2009complexity}. SVT identifies the first query in a sequence that exceeds a threshold. 
In the central model, a trusted curator adaptively screens intermediate outputs and halts upon finding a qualifying query, allowing SVT to consume only constant privacy budget regardless of sequence length. 
However, SVT is fundamentally incompatible with the shuffle model, which lacks a trusted curator. 
All intermediate query results must be exposed to the analyzer, causing privacy loss to accumulate linearly with the number of queries tested rather than remaining constant.


(2) \textit{Binary Search}~\cite{gilbert2002summarize, cormode2005improved}. 
An alternative approach performs binary search over the domain $[1, M]$ to approximate $m_{\max}(D)$. 
At each iteration, a noisy count query determines whether $m_{\max}(D)$ lies in the upper or lower half of the current search interval and refines the range accordingly. 
After $O(\log M)$ rounds, the algorithm identifies a threshold close to $m_{\max}(D)$ with high probability.
Such interactive protocols, however, incur prohibitive communication cost in shuffle-DP: each round requires all $n$ users to send messages, the shuffler to permute them, and the analyzer to compute intermediate results before proceeding to the next round.




\paragraph{Objective}
Therefore, our objective is to design a general shuffle-DP framework that privately estimates a threshold $m_\tau \approx m_{\max}(D)$ in constant rounds to achieve instance-specific error with limited communication cost.
% \section{Methodology}

% In this section, we first introduce the general framework for user-level shuffle-DP, beginning with the challenges and several naive solutions (but not actually work), and then introduce the core idea ...



\section{Framework for User-level Shuffle-DP}


This section addresses user-level shuffle-DP under the static setting.
We first present a two-round protocol (Section~\ref{sec: static two-round}) that achieves instance-specific error through interactive communication.
To align with the standard non-interactive assumption of shuffle-DP, we then propose a one-round protocol (Section~\ref{sec:static one-round}) that incurs only a logarithmic utility overhead.





\subsection{Two-round Protocol}
\label{sec: static two-round}
% please check that, in two-round protocol, we may not need domain partition.
% in two-round protocol, we can find the exact m_max(D) with probability 1-beta/2M
% no, bias will increase: loglog -> log; 
% so we may add a strawman which does not use domain partition.
% -> highlight challenge在哪，然后这个可以用domain partition解决，他为什么可以实现。



As discussed in Section~\ref{sec:Challenges and Objectives}, the core challenge is privately estimating an instance-specific $m_\tau$ that balances clipping bias and DP noise. We aim to satisfy:
\begin{enumerate}
    \item[(a)] The number of excluded records is bounded by $\tilde{O}(m_{\max}(D)/\varepsilon)$.
    \item[(b)] The DP noise scales with $\tilde{O}(m_{\max}(D)/\varepsilon)$.
\end{enumerate}
To achieve this, we adopt the \textit{domain partitioning} strategy, which has been successfully employed in shuffle-DP sum estimation~\cite{dong2024almost}.
Unlike binary search, which requires $O(\log M)$ rounds, domain partitioning exploits disjoint geometric subdomains to achieve single-round estimation with $O(1)$ messages per user.
The key idea is to partition $[1, M]$ into $\log M + 1$ geometrically increasing subdomains and evaluate user presence across all subdomains simultaneously.
% please check: emphasize record-level
Building on this strategy, we propose a two-round protocol with privacy budget $(\varepsilon/2, \delta/2)$ and failure probability $\beta/2$ per round to ensure overall $(\varepsilon,\delta)$-DP. 
The first round estimates $m_\tau$ via domain partitioning.
In the second round, users clip or pad their datasets to size $m_\tau$, after which we apply a base record-level shuffle-DP protocol $\mathcal{P}_Q$ with per-record privacy budget $(\varepsilon/2m_\tau, \delta/2m_\tau)$.



\subsubsection{Round $1$: Estimate $m_\tau$ via Domain Partitioning}
\label{sec: static round 1}
% \begin{small}
\begin{algorithm}[t]
\caption{Randomizer for Estimating $m_\tau$} 
\label{alg:Randomizer for Estimating tau}
% please check: this P is actually GKMPS. not a given mechanism.
% please check: eps or eps/2. how to write
\KwParam{$\varepsilon,\delta,n,M,\mathcal{P}_{\text{cnt}}$\cite{ghazi2021differentially}$=(\mathcal{R}_\mathrm{cnt},\mathcal{S}_\mathrm{cnt},\mathcal{A}_\mathrm{cnt})$}
\KwIn{$m_i$}
\For{$j \leftarrow 0$ \KwTo $\log M$}{
$C_i^{(j)} \leftarrow \mathcal{R}_\mathrm{cnt}(\mathbb{I}(m_i \in [2^{j-1}+1, 2^j]); \frac{\varepsilon}{2}, \frac{\delta}{2}, n)$;

\textbf{Send} $C_i^{(j)}$ to shuffler $\mathcal{S}_\text{cnt}^{(j)}$;
}
\end{algorithm}
% \end{small}

% \begin{small}
\begin{algorithm}[t]
\caption{Analyzer for Estimating $m_\tau$} 
\label{alg:Analyzer for Estimating tau}
\KwParam{$\varepsilon,\delta,\beta,n,M,\mathcal{P}_{\text{cnt}}$\cite{ghazi2021differentially}$=(\mathcal{R}_\mathrm{cnt},\mathcal{S}_\mathrm{cnt},\mathcal{A}_\mathrm{cnt})$}

\KwIn{$\{Z^{(j)} = \biguplus_{i=1}^n C_i^{(j)}\}_{j\in \{0,1,\dots,\log M\}}$}

$m_\tau \leftarrow 0$;

\For{$j \leftarrow 0$ \KwTo $\log M$}{

$\tilde{Q}_\mathrm{cnt}^{(j)} \leftarrow \mathcal{A}_\mathrm{cnt}(Z^{(j)}; \frac{\varepsilon}{2}, \frac{\delta}{2},\frac{\beta}{2(\log M +1)},n)$;

\If{$\tilde{Q}^{(j)}_\mathrm{count} > \frac{2}{\varepsilon}\ln\frac{2(\log M +1)}{\beta}$}{
$m_\tau \leftarrow 2^j$;
}

}

\Return $m_\tau$;
\end{algorithm}
% \end{small}

We partition $[1, M]$ into $\log M +1$ disjoint geometric subdomains\footnote{We assume $M$ is a power of 2. All $\log$ have base 2. Define $\log (0) := -1$ and $2^{-1}:= 0.$}:
$$I_j = [2^{j-1}+1, 2^j],\;j = 0,1,\dots,\log M.$$

\paragraph{Randomizer.} Each user $i$ computes the indicator $\mathbb{I}(m_i \in I_j)$ for each subdomain and reports it using the bit counting protocol $\mathcal{P}_{Q_\text{count}}$~\cite{ghazi2021differentially}. 
Since these subdomains are disjoint and each user contributes to exactly one, parallel composition allows allocating full budget $(\varepsilon/2, \delta/2)$ to each subdomain.
% please check, here need a description of shuffler.

\paragraph{Analyzer.} After receiving shuffled messages, the analyzer aggregates noisy counts for each subdomain, producing an estimate $\tilde{Q}^{(j)}_\text{count}$. 
We set $m_\tau = 2^j$ where $j$ is the maximum integer satisfying:
$$\tilde{Q}_\text{count}^{(j)} > \frac{2}{\varepsilon}\ln\frac{2(\log M+1)}{\beta}.$$
The right-hand side is derived from the error bound in Lemma~\ref{lem: Bit Counting Protocol}.
This threshold prevents overshooting---erroneously selecting empty subdomains---while ensuring that all unselected subdomains contain at most $\tilde{O}(1/\varepsilon)$ users; otherwise, their noisy counts would have exceeded the threshold.
The detailed algorithms for the randomizer and analyzer are presented in Algorithm~\ref{alg:Randomizer for Estimating tau} and Algorithm~\ref{alg:Analyzer for Estimating tau}.

Through this procedure, our protocol guarantees that $m_\tau$ satisfies:
\begin{enumerate}
\item[(a)] $\sum_{i=1}^n \max(0, m_i - m_\tau) \le \tilde{O}(m_{\max}(D)/\varepsilon)$
\item[(b)] $m_\tau \le 2 \cdot m_{\max}(D)$.
\end{enumerate}
Detailed proofs are in Appendix~\ref{sec:Proof of static two}.



% Intuitively, $m_\tau$ serves as a private upper bound on $m_{\max}(D)$: with high probability,
% no user has more than $m_\tau$ records.
% please check that here may need a analysis on the m_tau. we need to prove that the tau is suitable. e.g. loglog(M). may refer to CCS.

% \begin{small}
\begin{algorithm}[t]
\caption{Randomizer for Query $Q$} 
\label{alg: randomizer for query evaluation}
\KwParam{$\varepsilon,\delta,n,\mathcal{P}_{Q}=(\mathcal{R}_Q,\mathcal{S}_Q,\mathcal{A}_Q)$}
\KwIn{$m_\tau$, $D_i = (x_{i,1}, x_{i,2}, \ldots, x_{i,m_i})$}


\eIf{$m_i < m_\tau$}{
    $D_i \gets D_i \mathbin{\|} [\bot]^{(m_\tau - m_i)}$ ; 
}{
    $D_i \gets D_i[1:m_\tau]$ ;
}


\For{$k \gets 1$ \KwTo $m_\tau$}{
    $Y_{i,k} \gets \mathcal{R}_Q(D_i[k]; \,\frac{\varepsilon}{2m_\tau},\, \frac{\delta}{2m_\tau},\, n\cdot m_\tau)$;
    
    \textbf{Send} $Y_{i,k}$ to shuffler $\mathcal{S}_\text{qry}$;
}

\end{algorithm}
% \end{small}


% \begin{small}
\begin{algorithm}[t]
\caption{Analyzer for Query $Q$} 
\label{alg:analyzer for query evaluation}
\KwParam{$\varepsilon,\delta,\beta,n,M,\mathcal{P}_{Q}=(\mathcal{R}_Q,\mathcal{S}_Q,\mathcal{A}_Q)$}

\KwIn{$Z = \biguplus_{i,k} Y_{i,k}$}

$\tilde{Q}\gets \mathcal{A}_Q(Z; \frac{\varepsilon}{2m_\tau},\, \frac{\delta}{2m_\tau},\, \frac{\beta}{2},  n\cdot m_\tau)$;

\Return $\tilde{Q}$

\end{algorithm}
% \end{small}

% please check, i think the algorithms here maybe no need.


\subsubsection{Round $2$: Contribution Standardization and Query Evaluation}


With $m_\tau$ as the input parameter, we reduce the user-level DP problem to a record-level one by standardizing each user's contribution to exactly $m_\tau$ records. Specifically, for each user $i$, the standardized dataset $D_i^\tau$ is defined as:
$$D^\tau_i =
\begin{cases}
D_i[1:m_\tau] & \text{if } m_i > m_\tau, \\
D_i \mathbin{\|} (\bot)^{m_\tau - m_i} & \text{if } m_i \le m_\tau.
\end{cases}
$$
where $\|$ denotes the concatenation operation. The resulting dataset is $D^\tau = \bigcup_i D^\tau_i$. For any pair of user-level neighbors $D \sim_u D'$, we have $d(D^{\tau}, D'^\tau) \le m_\tau$. Consequently, by group privacy (Lemma~\ref{lem:Group Privacy}), we can apply any record-level shuffle-DP protocol $\mathcal{P}_Q$ on $D^\tau$ with privacy budget $(\frac{\varepsilon}{2m_\tau}, \frac{\delta}{2m_\tau})$. The detailed algorithms for the randomizer and analyzer are presented in Algorithm~\ref{alg: randomizer for query evaluation} and Algorithm~\ref{alg:analyzer for query evaluation}.




%please check: does binary search need to split eps by logM?

In summary, our two-round protocol addresses the user-level shuffle-DP problem, achieving instance-specific error that depends on $m_{\max}(D)$ rather than the global bound $M$. Moreover, it maintains efficient communication: by exploiting the disjoint structure of geometric subdomains in Round 1, each user sends only $O(1)$ messages to estimate $m_\tau$.



\begin{theorem}
\label{the: static two-round}   
    Given any $\varepsilon > 0, \delta >0, n \in \mathbb{Z}_+, $ $U \in \mathbb{Z}_{+}$, $M \in \mathbb{Z}_{+}$, for any query $Q$ and any base shuffle-DP protocol $\mathcal{P}_Q$, our two-round protocol satisfies the following properties in the static setting:

    \begin{itemize}
    \item The messages received by the analyzer preserve $(\varepsilon,\delta)$-DP;
    \item The total error of this protocol is bounded by:
    \[
    \begin{aligned}
    &\mathrm{\gamma}_{\ell_p}\!\Big(
        Q,\;
        \frac{16}{\varepsilon}\ln\frac{2(\log M +1)}{\beta} \cdot m_{\max}(D)
    \Big)
    \\
    &+\;
    \err_{\ell_p}\!\Big(
        \mathcal{P}_Q,\;
        \frac{\varepsilon}{4m_{\max}(D)},\;
        \frac{\delta}{4m_{\max}(D)},\;
        \frac{\beta}{2},\;
        n\cdot 2m_{\max}(D)
    \Big)
    \end{aligned}
    \]



    \item In expectation, each user sends $1+ o(1)+\mathrm{Msg}(\mathcal{P}_Q, \allowbreak \frac{\varepsilon}{4m_{\max}(D)}, \allowbreak \frac{\delta}{4m_{\max}(D)}, \allowbreak n\cdot 2m_{\max}(D))$ messages.

\end{itemize}

\end{theorem}


Due to space limitations, all proofs are deferred to the Appendix.


\begin{algorithm}[t]
\caption{Randomizer}
\label{alg:Randomizer for one round}
\KwParam{$\varepsilon,\delta,n,M,\mathcal{P}_{Q_\text{cnt}}$\cite{ghazi2021differentially}$=(\mathcal{R}_\mathrm{cnt},\mathcal{S}_\mathrm{cnt},\mathcal{A}_\mathrm{cnt})$} 
\KwIn{ $D_i = (x_{i,1}, x_{i,2}, \ldots, x_{i,m_i})$, $\mathcal{P}_{Q}=(\mathcal{R}_Q,\mathcal{S}_Q,\mathcal{A}_Q)$}

\For{$j \leftarrow 0$ \KwTo $\log M$}{
$C_i^{(j)} \leftarrow \mathcal{R}_{\text{count}}(\mathbb{I}(m_i \in [2^{j-1}+1, 2^j]); \frac{\varepsilon}{2},\frac{\delta}{2}, n)$;


\textbf{Send} $C_i^{(j)}$ to shuffler $\mathcal{S}_\text{cnt}^{(j)}$;


\eIf{$m_i < 2^j$}{
    $\tilde{D} \gets D_i \mathbin{\|} [\bot]^{(2^j - m_i)}$ ; 
}{
    $\tilde{D} \gets D_i[1:2^j]$ ;
}





\For{$k \gets 1$ \KwTo $2^j$}{
    $Y^{(j)}_{i,k} \gets \mathcal{R}_Q(\tilde{D}[k]; \,\frac{\varepsilon}{2\cdot2^j \cdot (\log M+1) },\, \frac{\delta}{2 \cdot2^j \cdot (\log M+1) },\, n\cdot 2^j)$;
    
    {\textbf{Send}} $Y^{(j)}_{i,k}$ to shuffler $\mathcal{S}_\text{qry}^{(j)}$;
}




}

\end{algorithm}




\subsection{A One-round Protocol}
\label{sec:static one-round}
While the two-round protocol effectively solves the user-level shuffle-DP problem, it requires interactive communication in which the analyzer sends $m_\tau$ back to users.
However, the standard shuffle-DP model assumes a one-way communication channel from users to the shuffler and then to the analyzer.
% Please check, here may need a citation.
To adhere to this non-interactive assumption, we propose a one-round framework.
% where users compute query responses for all $\log M$ candidate subdomains simultaneously.
% The analyzer then identifies $m_\tau$ and selectively aggregates corresponding query responses.
Although this approach requires splitting the privacy budget across subdomains, we show that the additional utility overhead is only a logarithmic factor.



\paragraph{Randomizer.} 
Each user splits the privacy budget $(\varepsilon, \delta)$ equally between two tasks: 
(1) estimating $m_\tau$ using the same domain partitioning strategy as in Round 1 of the two-round protocol, and 
(2) preparing query responses for all candidate thresholds simultaneously.
% The key difference from the two-round protocol is that users must prepare responses for all $\log M + 1 $candidate thresholds simultaneously without knowing which threshold the analyzer will select.
Since the query evaluations overlap across subdomains---the dataset prepared for threshold $2^{j+1}$ subsumes that for $2^j$---basic composition requires dividing the query evaluation budget equally among $\log M+1$ subdomains.
For each subdomain $I_j$, the user constructs a temporary dataset $\tilde{D}$ by truncating or padding $D_i$ to size $2^j$, then applies $\mathcal{R}_Q$ with privacy budget $(\frac{\varepsilon}{2 \cdot 2^j \cdot (\log M + 1)},\frac{\delta}{2 \cdot 2^j \cdot (\log M + 1)})$.
The detailed procedure is shown in Algorithm~\ref{alg:Randomizer for one round}.

% \begin{small}
\begin{algorithm}[t]
\caption{Analyzer} 
\label{alg:Analyzer for One round}
\KwParam{$\varepsilon,\delta,\beta,n,M,\mathcal{P}_{Q_\text{cnt}}$\cite{ghazi2021differentially}$=(\mathcal{R}_\mathrm{cnt},\mathcal{S}_\mathrm{cnt},\mathcal{A}_\mathrm{cnt})$}

\KwIn{$\{Z^{(j)}_\text{count} = \biguplus_{i=1}^n C_i^{(j)}\}_{j\in \{0,1,\dots,\log M\}}$, $\{Z^{(j)}_Q = \biguplus_{i,k} Y^{(j)}_{i,k}\}_{j\in \{0,1,\dots,\log M\}}$, $\mathcal{P}_{Q}=(\mathcal{R}_Q,\mathcal{S}_Q,\mathcal{A}_Q)$}

$\beta' \leftarrow \frac{\beta}{2(\log M+1)}$;

$m_\tau \leftarrow 0$;

\For{$j \leftarrow 0$ \KwTo $\log M$}{

$\tilde{Q}^{(j)}_\text{cnt} \leftarrow \mathcal{A}_\text{cnt}(Z_\text{cnt}^{(j)}; \frac{\varepsilon}{2},\frac{\delta}{2},\beta',n)$;

\If{$\tilde{Q}^{(j)}_\mathrm{cnt} > \frac{2}{\varepsilon}\ln\frac{1}{\beta'}$}{
$m_\tau \leftarrow 2^j$;
}

}

$\tilde{Q} \gets \mathcal{A}_Q(Z_Q^{(\log m_\tau)}; \frac{\varepsilon}{2 (\log M+1) \cdot m_\tau},\, \frac{\delta}{2 (\log M+1) \cdot m_\tau},\beta',  n\cdot m_\tau)$;

\Return $\tilde{Q}$

\end{algorithm}
% \end{small}


\paragraph{Analyzer.}
The analyzer's procedure is presented in Algorithm~\ref{alg:Analyzer for One round}.
Upon receiving the shuffled messages $\{Z^{(j)}_\text{count}\}$ and $\{Z^{(j)}_Q\}$, the analyzer proceeds in two steps:
\begin{enumerate}
\item \textit{Determining $m_\tau$.} 
The analyzer first invokes $\mathcal{A}_\text{count}$~\cite{ghazi2021differentially} on $Z^{(j)}_\text{count}$ for each subdomain $I_j$ to obtain the estimated user count $\tilde{Q}^{(j)}_\text{count}$. 
Following the same strategy as the two-round protocol, it identifies the largest index $j$ where the estimated count exceeds $\frac{2}{\varepsilon}\ln\frac{2 \cdot (\log M+1)}{\beta}$. 
% please check here should be query evaluation or selective aggregation
\item \textit{Selective Aggregation.} Once $m_\tau$ is determined, the analyzer discards query messages from all other subdomains and aggregates only $Z_Q^{(\log m_\tau)}$ to compute the final estimate $\tilde{Q}$.
\end{enumerate}


\begin{theorem}
\label{the: static one-round}   
    Given any $\varepsilon > 0, \delta >0, n \in \mathbb{Z}_+, $ $U \in \mathbb{Z}_{+}$, $M \in \mathbb{Z}_{+}$, under the static setting, our one-round protocol achieves the following guarantees:

    \begin{itemize}
    \item The messages received by the analyzer preserve $(\varepsilon,\delta)$-DP;
    \item The total error of this protocol is bounded by:
    \[
    \begin{aligned}
    &\mathrm{\gamma}_{\ell_p}\!\Big(
        Q,\;
        \frac{16}{\varepsilon}\ln\frac{2(\log M+1)}{\beta} \cdot m_{\max}(D)
    \Big)
    +\;
    O\Big(\err_{\ell_p}\!(
        \mathcal{P}_Q,\;    \\
    &
        \frac{\varepsilon}{m_{\max}(D)\cdot \log M},\;
        \frac{\delta}{m_{\max}(D)\cdot \log M},\;
        \frac{\beta}{\log M},\;
        n\cdot m_{\max}(D)
    )\Big);
    \end{aligned}
    \]



    \item In expectation, each user sends $1+o(1)+\sum_{j \in \{0,1,\dots, \log M\}} \allowbreak \mathrm{Msg}(\mathcal{P}_Q, \frac{\varepsilon}{2\cdot 2^j(\log M+1)}, \frac{\delta}{2\cdot 2^j(\log M+1)},n \cdot 2^j)$ messages.

\end{itemize}

\end{theorem}



\section{Framework for Dynamic User-level Shuffle-DP}

\label{sec:Framework for Dynamic User-level Shuffle-DP}


In this section, we look at the dynamic setting. 
This setting introduces a fundamental challenge: continuously adapt $m_\tau^t$ at each timestep to track the time-varying $m_{\max}(D^t)$.
We begin by reviewing how central-DP addresses this challenge, and then explain why it fails in the shuffle-DP model, and finally present our protocol.



\paragraph{Central-DP Solution}
\label{sec:dynamic centralDP solution}

Central-DP uses SVT~\cite{dong2023continual} to track $m_{\max}(D^t)$. 
The curator continuously checks whether $m_{\max}(D^t)$ exceeds the current threshold $\tau_t$. 
Upon detecting a crossing, it updates $m_\tau^t \leftarrow 2\tau_t$ and launches a new SVT instance with reduced budget $\varepsilon_i = \varepsilon\theta/(i+1)^{1+\theta}$. 
Crucially, SVT consumes privacy budget only when crossings occur, not at every timestep, enabling efficient long-horizon tracking.
However, as discussed in Section~\ref{sec:Challenges and Objectives}, SVT cannot be adapted to shuffle-DP. 
Without a trusted curator, all intermediate noisy counts must be revealed to the analyzer, causing privacy loss at every timestep and rendering the approach impractical.



\subsection{Our Protocol}



Our key insight is to decompose adaptive $m_\tau$ estimation into concurrent dynamic counting queries.
Recall that, in our static protocol (Section~\ref{sec:static one-round}), we partitioned users into $\log M+1$ buckets corresponding to contribution subdomains $I_j = [2^{j-1}+1, 2^j]$ for $j \in \{0, 1, \dots, \log M\}$, computing one fixed count per bucket. 
Using the binary mechanism~\cite{ghazi2021differentially} (Section~\ref{sec:Binary mechanism}), we extend these static buckets into $\log M+1$ dynamic counting streams. 
By maintaining $\log M+1$ parallel binary tree instances, we continuously track user distributions across subdomains at any timestamp $t$.


Building upon this insight, similar as our one-round framework under static setting, our framework executes two tasks simultaneously: 
\textit{estimating $m_\tau^t$} by transforming the $\log M + 1$ static counting buckets into dynamic counting streams, and \textit{query evaluation} by continuously evaluating the query at all $\log M + 1$ candidate thresholds using the dynamic event-level shuffle-DP protocol.
This structure enables the analyzer to estimate $m_\tau$ at any timestamp and retrieve the corresponding query result.







%please check: think the name
\begin{algorithm}[t]
\caption{Dynamic Counting (User Side)}
\label{alg:dynamic_counting}
\KwIn{$\varepsilon, \delta, j, n, D_i = \{x_{i,t}\}_{t=1}^T, \mathcal{P}_\mathrm{cnt}$}
\SetKwData{Acc}{Acc}
\SetKwFunction{Bit}{Bit}
\SetKwFunction{RCnt}{$\mathcal{R}_\mathrm{cnt}$}

$\Acc[h] \leftarrow 0$, $\forall h \in \{0, 1, \dots, \log T\}$;\enspace $m_i \leftarrow 0$;\enspace $j_{\text{curr}} \leftarrow 0$;

\For{$t \leftarrow 1$ \KwTo $T$}{
    $j_{\text{pre}} \leftarrow j_{\text{curr}}$;
    
    \If{$x_{i,t} \neq \perp$}{
        $m_i \leftarrow m_i + 1$;
        
        $j_{\text{curr}} \leftarrow \lfloor \log m_i \rfloor + 1$;
        
        $\Acc[h] \leftarrow \Acc[h] + \mathbb{I}[j_{\text{curr}}=j] - \mathbb{I}[j_{\text{pre}}=j]$,\enspace $\forall h$;
    }
    
    $h^* \leftarrow \min \{ h : \Bit_h(t) = 1 \}$;

    % $C_i^{(j,t)} \leftarrow \RCnt\!\left(\Acc[h^*], \frac{\varepsilon}{\log T+1}, \frac{\delta}{\log T+1}, n\right)$;
    
    % \textbf{Send} $C_i^{(j,t)}$ to shuffler $\mathcal{S}^{(j)}_{\text{cnt}}$;
    
    \textbf{Send} $\RCnt\big(\Acc[h^*], \frac{\varepsilon}{\log T+1}, \frac{\delta}{\log T+1}, n\big)$ to shuffler $\mathcal{S}^{(j)}_{\text{cnt}}$;
    
    $\Acc[0:h^*] \leftarrow \mathbf{0}$;
}
\end{algorithm}


% please check, 统一一下[]的用法，[]_0

\begin{algorithm}[t]
\caption{Query Evaluation (User Side)}
\label{alg:query_evaluation}
\KwIn{$\varepsilon, \delta, j, n, D_i = \{x_{i,t}\}_{t=1}^T,\mathcal{P}_Q$}
\SetKwData{Acc}{Acc}
\SetKwFunction{Bit}{Bit}
\SetKwFunction{RQ}{$\mathcal{R}_Q$}


$m_i \leftarrow 0$; $m_\tau=2^j$;

$S[k] \leftarrow (\perp)^T$, $\forall k \in \{1,2,\dots,m_\tau\}$;

\textbf{Start} \RQ$(S[k], \frac{\varepsilon}{m_\tau}, \frac{\delta}{m_\tau}, n\cdot m_\tau )$ to $\mathcal{S}^{(j)}_{\mathrm{qry}}$, $\forall  k$;


\For{$t \leftarrow 1$ \KwTo $T$}{

    \If{$x_{i,t} \neq \perp $ \KwAnd $ m_i < m_\tau$}{
        $m_i \leftarrow m_i + 1$;
        
        $S[m_i][t] \leftarrow x_{i,t}$;
    }
    
    \textbf{Continue} \RQ$(S[k], \frac{\varepsilon}{m_\tau}, \frac{\delta}{m_\tau}, n\cdot m_\tau )$, $\forall k$;
}

\end{algorithm}

\begin{algorithm}[t]
\caption{Randomizer}
\label{alg:randomizer}
\KwIn{$\varepsilon, \delta, T, M, n, D_i = \{x_{i,t}\}_{t=1}^T, \mathcal{P}_\mathrm{cnt},\mathcal{P}_Q$}
\SetKwFunction{DynCount}{DynamicCounting}
\SetKwFunction{QueryEva}{QueryEvaluation}
\SetKwFunction{RQ}{$\mathcal{R}_Q$}

$\varepsilon' \leftarrow \frac{\varepsilon}{2(\log M+1)}$, $\delta' \leftarrow\frac{\delta}{2(\log M+1)}$;


\textbf{Start} \DynCount$(D_i,\varepsilon', \delta', j, n)$, $\forall j \in [\log M]$;


\textbf{Start} \QueryEva$(D_i, \varepsilon', \delta', j, n)$, $\forall j$;

\end{algorithm}



\paragraph{Randomizer}



The randomizer allocates equal privacy budget $(\varepsilon/2, \delta/2)$ to each task. 
Since a user may participate across all $\log M+1$ subdomains, each task requires sequential composition, yielding per-subdomain budget $(\varepsilon', \delta') = \big(\frac{\varepsilon}{2(\log M+1)}, \frac{\delta}{2(\log M+1)}\big)$. 
Algorithm~\ref{alg:randomizer} shows how the two tasks are instantiated across all subdomains.
The sub-procedures for each subdomain $j$ are given in Algorithm~\ref{alg:dynamic_counting} (Dynamic Counting) and Algorithm~\ref{alg:query_evaluation} (Query Evaluation).


\textit{(1) Dynamic Counting.}
At each timestep $t$, the user sets $j_{\text{pre}} \leftarrow j_{\text{curr}}$. 
When new data $x_{i,t} \neq \perp$ arrives, the user increments $m_i$, updates $j_{\text{curr}} = \lfloor \log m_i \rfloor + 1$, then accumulates the update $\mathbb{I}[j_{\text{curr}}=j] - \mathbb{I}[j_{\text{pre}}=j]$ into all $\log T+1$ binary-tree levels, where $+1$ indicates entering $I_j$ and $-1$ indicates leaving. 
Each tree level $h$ corresponds to time windows of length $2^h$, enabling count reconstruction at any timestamp.
At each $t$, the user identifies flush node $h^* = \min\{h : \text{Bit}_h(t) = 1\}$, where $\text{Bit}_h(t)$ denotes the $(h+1)$-th least significant bit of $t$,\footnote{For example, at $t=6$ (binary: $110_2$), the rightmost $1$ appears at position $h^*=1$, triggering an update covering time interval $\{5,6\}$.} privatizes $\text{Acc}[h^*]$ via $\mathcal{R}_{\text{cnt}}$ with pricavy budget $(\frac{\varepsilon'}{\log T+1}, \frac{\delta'}{\log T+1})$, sends it to $\mathcal{S}^{(j)}_{\text{cnt}}$, and resets $\text{Acc}[0:h^*] \leftarrow \mathbf{0}$, preparing to accumulate counts for next time window.


\textit{(2) Query Evaluation.}
For each candidate threshold $m_\tau = 2^j$, user $i$ maintains $m_\tau$ streams $\{S[k]\}_{k=1}^{m_\tau}$, each a length-$T$ sequence initially filled with $\perp$. 
The user then starts the base event-level shuffle-DP protocol instances $\mathcal{R}_Q$ for all streams, each with privacy budget $(\frac{\varepsilon'}{m_\tau}, \frac{\delta'}{m_\tau})$ and participant size $n \cdot m_\tau$. 
When data $x_{i,t} \neq \perp$ arrives and $m_i \le m_\tau$, the user assigns the new data to the $m_i$-th stream: $$S[m_i][t] \leftarrow x_{i,t};$$ otherwise the data is discarded.
After that, all $m_\tau$ streams continue running the base protocol.\footnote{This incurs no additional privacy cost because each position is written at most once: $S[m_i][t]$ transitions from $\perp$ to $x_{i,t}$ exactly when the datum first arrives, before participating in any  shuffle-DP mechanism.}


\begin{algorithm}[t]
\caption{Analyzer}
\label{alg:analyzer}
\KwParam{$\varepsilon,\delta,\beta,T, M, n$}
\KwIn{$\{{Z}_{\mathrm{cnt}}^{(j,t)}\}_{j,t}$ from shuffler $\mathcal{S}^{(j)}_{\mathrm{cnt}}$; \\
$\{{Z}_{\mathrm{qry}}^{(j,t)}\}_{j,t}$ from shuffler $\mathcal{S}^{(j)}_{\mathrm{qry}}$
}
\SetKwFunction{Acnt}{$\mathcal{A}_\mathrm{cnt}$}
\SetKwFunction{Bit}{Bit}

% Initialization
$\hat{C}[j][h] \leftarrow 0$ for all $j \in \{0,\dots,\log M\}, h \in \{0,\dots,\log T\}$\;

$\varepsilon' \leftarrow \frac{\varepsilon}{2(\log M+1)}$, $\delta' \leftarrow\frac{\delta}{2(\log M+1)}$, $\beta' \leftarrow \frac{\beta}{2(\log M+1)}$;

\textbf{Start} $\mathcal{A}_Q^{(j)}(\{{Z}_{\mathrm{qry}}^{(j,t)}\}_t,\frac{ \varepsilon'}{2^j}, \frac{\delta'}{2^j}, \beta', n\cdot 2^j)$ for each $j$;

\For{$t \leftarrow 1$ \KwTo $T$}{
    \tcp{Task 1: Reconstruct interval counts}
    $h^* \leftarrow \min \{ h : \Bit_h(t) = 1 \}$;
    
    \For{$j \leftarrow 0$ \KwTo $\log M$}{

        $\hat{C}[j][h^*] \leftarrow \Acnt({Z}_{\mathrm{cnt}}^{(j,t)}, \frac{\varepsilon'}{\log T+1}, \frac{\delta'}{\log T+1}, \frac{\beta'}{T}, n)$;

        $\tilde{Q}_\mathrm{cnt}^{(j)}  \leftarrow \sum_{h: \Bit_h(t)=1} \hat{C}[j][h]$;
    }

    $j^* \leftarrow \max \big\{0,  j : \tilde{Q}_\mathrm{cnt}^{(j)} > \frac{\log ^{1.5} T}{\varepsilon'}\ln \frac{T}{\beta'} \big\}$;
    % please check here, this term is $O(..)$
    
    \tcp{Task 2: Choose query result for $m_\tau=2^{j^*}$}
    \textbf{Get} $\tilde{Q}$ from $ \mathcal{A}_Q^{(j^*)}$;
    
    \textbf{Output} $\tilde{Q}$;
}
\end{algorithm}




\paragraph{Analyzer}
The analyzer receives two sets of shuffled messages at each time $t$: $Z^{(j,t)}_\mathrm{cnt}$ from shuffler $\mathcal{S}^{(j)}_\mathrm{cnt}$ for estimating $m_\tau$, and $Z^{(j,t)}_\mathrm{qry}$ from shuffler $\mathcal{S}^{(j)}_\mathrm{qry}$ for query evaluation.
To identify $m_\tau^t$, the analyzer maintains a matrix $\hat{C}[j][h]$ storing the noisy count for interval $j$ at tree level $h$. 
At each $t$, it identifies the flush node $h^* = \min\{h : \mathrm{Bit}_h(t) = 1\}$, updates the stored count:
$$\hat{C}[j][h^*] \leftarrow \mathcal{A}_{\mathrm{cnt}}\!\left(Z_{\mathrm{cnt}}^{(j,t)},\, \tfrac{\varepsilon'}{\log T+1},\, \tfrac{\delta'}{\log T+1},\, \tfrac{\beta'}{T},\, n\right),$$
and recovers the current count for each $j$ by binary tree aggregation:
$$\tilde{Q}_{\mathrm{cnt}}^{(j)} = \sum_{h:\,\mathrm{Bit}_h(t)=1} \hat{C}[j][h].$$
Here $\tilde{Q}_{\mathrm{cnt}}^{(j)}$ estimates the number of users with total contribution $m_i$ falls in $I_j$ at time $t$. 
The analyzer selects $$j^* = \max\!\left\{0,\; j : \tilde{Q}_{\mathrm{cnt}}^{(j)} > \frac{\log^{1.5} T}{\varepsilon'}\ln\frac{T}{\beta'}\right\},$$
which identifies $m_{\max}(D^t) \in I_{j^*}$ with probability at least $1 - \beta'/2$.
Since all $\log M + 1$ Query Evaluation instances run continuously simultaneously, the query result under threshold $2^{j^*}$ is already available: the analyzer simply retrieves $\tilde{Q}$ from $\mathcal{A}_Q^{(j^*)}$ and outputs it.
The detailed procedure is shown in Algorithm~\ref{alg:analyzer}.






\begin{theorem}
\label{the: dynamic setting guarantee}
    Given any $\varepsilon > 0, \delta > 0$, any $T \in \mathbb{Z}_+$, and $n \in \mathbb{Z}_+$, $M \in \mathbb{Z}_+$, for any query $Q$, at all timestamps $t \in \{1, 2, \dots, T\}$ simultaneously, our protocol for the dynamic setting achieves the following guarantees:
\begin{itemize}
    \item The messages received by the analyzer preserve $(\varepsilon,\delta)$-DP;
    \item With probability at least $1-\beta$, the total error is bounded by:
    \[
    \begin{aligned}
    &O\Big(\gamma_{\ell_p}\big( Q, \frac{\log ^{1.5} T \cdot \log M}{\varepsilon} \ln \frac{T \cdot \log M}{\beta} \cdot m_{\max}(D^t)\big)\Big) + O\Big(\err_{\ell_p}\big(\mathcal{P}_Q,
    \\
    & \frac{\varepsilon}{m_{\max}(D^t) \cdot \log M}, \frac{\delta}{m_{\max}(D^t) \cdot \log M}, \frac{\beta}{\log M}, n \cdot m_{\max}(D^t)\big)\Big)
    \end{aligned}
    \]
    \item In expectation, each user sends $(\log M+1)(1+o(1)) + \sum_{j=0}^{\log M} \sum_{k=1}^{2^j} \mathrm{Msg}(\mathcal{P}_Q, \frac{\varepsilon}{2\cdot 2^j(\log M+1)}, \frac{\delta}{2\cdot 2^j(\log M+1)}, n \cdot 2^j)$ messages over the entire time horizon $T$.
\end{itemize}
\end{theorem}


% - range counting：Ghazi - On the Power of Multiple Anonymous Messages
%     - 构建了一颗range 二叉树，每个节点代表一个range，每个节点都做一次shuffle-DP，任意range可以由这颗二叉树上的O(log)个节点恢复
%     - 如何填dummy message：可以直接用dummy符号，或者2U。
% - median：
%     - 二分：Instance-optimal Mean Estimation Under Differential
%     Privacy


\section{Application}
In this section, we apply our framework for static setting in three queries: frequency estimation, range counting, and median, as defined in Section~\ref{sec:Notation}.
For dynamic setting, we apply our framework in bit counting query. We assume that $\varepsilon = \Theta(1)$ and $\log n = \Theta(\log(1/\delta)) = \Theta(\log M)= \Theta(\log T)$.



\subsection{Static Setting}
% please check: which framework should be used to apply to this 3 queries? multiple or single round. 
% can write both


\paragraph{Frequency estimation}
% 1) run B times sum estimation
% 2) qiyao's fe1
% pure-dp use 1; unless use 2
% we do not care pure/approximate DP, so fe1


For the frequency estimation problem, the SOTA solution LWY \cite{luo2022frequency} sends $O(1)$ messages and achieves an $\ell_{\infty}$-error of
$O\big(\tfrac{1}{\varepsilon}\sqrt{\log \tfrac{U}{\beta} \log \tfrac{1}{\delta}}\big)$.
By integrating it into our two-round framework, we have:
\begin{itemize}
% please check the bias
% read the other paper on how to represent the bias
\item With high probability $1-\beta$, the total error is bounded by: 
$$
% O\Big(\frac{m_{\max}(D)}{\varepsilon}\cdot \big(\log\frac{\log M}{\beta}+\sqrt{\log \frac{U}{\beta}\log\frac{m_{\max}(D)}{\delta}}\big)\Big).
O\Big(\frac{m_{\max}(D)}{\varepsilon}\cdot \big(\log\frac{\log M}{\beta}+\sqrt{\log \frac{U}{\beta}\log\frac{1}{\delta}}\big)\Big).
$$
\item The expected number of messages is $O(1)$.
\end{itemize}

By integrating it into our one-round framework, we have:

\begin{itemize}
\item With high probability $1-\beta$, the total error is bounded by: 
$$
% O\Big(\frac{m_{\max}(D)}{\varepsilon}\big(\log \frac{1}{\beta}+\log M \cdot \sqrt{
% \log \frac{U\cdot\log M }{\beta}
% \cdot \log \frac{\log M \cdot m_{\max}(D)}{\delta} }\big)\Big).
O\Big(\frac{m_{\max}(D)}{\varepsilon}\big(\log \frac{1}{\beta}+\log M \cdot \sqrt{
\log \frac{U\cdot\log M }{\beta}
\cdot \log \frac{1}{\delta} }\big)\Big).
$$
% $$O\Big(\frac{m_{\max}(D)}{\varepsilon}\big(\log \frac{\log M}{\beta}+\log M \cdot \sqrt{
% \log \frac{U\cdot\log M }{\beta}
% \cdot \log \frac{\log M \cdot m_{\max}(D)}{\delta} }\big)\Big).$$

\item The expected number of messages is $O(\log M)$.
\end{itemize}


\paragraph{Range counting}
% on the power of Ghazi'20
% rm2 better, which is an advanced version of <on the power ...>

Range counting can be reduced to frequency estimation by building a binary tree over the domain $[0,U]$, where each node represents a range and each level forms a frequency estimation problem.
The SOTA shuffle-DP protocol LYD+ \cite{luo2025rm2} achieves an $\ell_{\infty}$-error of
$O\big(\tfrac{\log^{1.5} U}{\varepsilon} \sqrt{\log \tfrac{U}{\beta} \log \tfrac{1}{\delta}}\big)$, with $O(\log U)$ messages per user.
By integrating it into our two-round framework, we have:
\begin{itemize}
% please check the bias
% read the other paper on how to represent the bias
\item With high probability $1-\beta$, the total error is bounded by: 
$$
% O\Big(\frac{m_{\max}(D)}{\varepsilon}\cdot \big(\log\frac{\log M}{\beta}+\log^{1.5}U\cdot\sqrt{\log \frac{U}{\beta}\log\frac{m_{\max}(D)}{\delta}}\big)\Big).
O\Big(\frac{m_{\max}(D)}{\varepsilon}\cdot \big(\log\frac{\log M}{\beta}+\log^{1.5}U\cdot\sqrt{\log \frac{U}{\beta}\log\frac{1}{\delta}}\big)\Big).
$$
\item The expected number of messages is $O(\log U)$.
\end{itemize}

By integrating it into our one-round framework, we have:
% 可以把多个项分开，我们的framework只引入了一个更小的可以被dominate的term
\begin{itemize}
\item With high probability $1-\beta$, the total error is bounded by: 
% \[
% \begin{aligned}
% &O\Bigg(
% \frac{m_{\max}(D)}{\varepsilon}\Big(
% \log \frac{\log M}{\beta} + \log^{1.5} U \cdot \log M \cdot \\
% &\sqrt{
% \log \frac{U\cdot\log M}{\beta}
% \cdot \log \frac{\log M \cdot m_{\max}(D)}{\delta}
% }
% \Big)
% \Bigg)
% \end{aligned}
% \]
\[
\begin{aligned}
O\Bigg(
\frac{m_{\max}(D)}{\varepsilon}\Big(
\log \frac{1}{\beta} + \log^{1.5} U \cdot \log M \cdot \sqrt{
\log \frac{U\cdot\log M}{\beta}
\cdot \log \frac{1}{\delta}
}
\Big)
\Bigg)
\end{aligned}
\]
% $$O\Big(\frac{m_{\max}(D)}{\varepsilon}\big(\log \frac{\log M}{\beta}+\log M \cdot \sqrt{
% \log \frac{U\cdot\log M }{\beta}
% \cdot \log \frac{\log M \cdot m_{\max}(D)}{\delta} }\big)\Big).$$

\item The expected number of messages is $O(\log U \log M)$.
\end{itemize}


\paragraph{Median}
% why median not mean
% instance-optimal mean estimation under differential privacy. Ziyue
% use <on the power of ...> to execute frequency estimation, in a binary structure.
% but multiple round

Median can be reduced to range counting by finding a value $v$ such that the number of elements smaller than $v$ is approximately half of the total. The SOTA shuffle-DP protocol HLY~\cite{huang2021instance} achieves a rank error of $\frac{\log^{2.5} U}{\varepsilon} \sqrt{\log \frac{U}{\beta}\log \frac{1}{\delta}}$, with $O(\log ^5 U \log \frac{1}{\delta})$.
% here we assume O(log U) = O(log 1/\delta)
By integrating it into our two-round framework, we have:
\begin{itemize}
\item With high probability $1-\beta$, the rank error is bounded by: $$O\Big(\frac{m_{\max}(D)}{\varepsilon}\big( \log \frac{\log M}{\beta}+\log^3 U \sqrt{\log \frac{U\log M}{\beta}\log \frac{1}{\delta}}\big)\Big).$$
\item The expected number of messages is $O(\log ^5 U \log \frac{1}{\delta})$.
\end{itemize}

% use rank error or inf error?
% log times range counting

\subsection{Dynamic Setting}
In the dynamic setting, we focus on the bit counting query.
The SOTA shuffle-DP protocol for bit counting in the dynamic setting TKMS~\cite{pmlr-v202-tenenbaum23a} achieves an error 
% $O(\frac{\sqrt{\log (1/\beta)}\log^{1.5}T+\log^2T}{\varepsilon}\sqrt{\log \frac{1}{\delta}})$
$O(\frac{\log^{1.5}T}{\varepsilon}\sqrt{\log \frac{1}{\delta}})$
with $O(t\log T)$ messages\footnote{Unlike the general dynamic setting, TKMS assumes a public user-to-time mapping where user $i$ contributes exactly once at time $t=i$. 
This is easily adapted to general setting: let all users participate at every time step $t\in[T]$, with user $i$ reporting their actual bit when $t=i$ and a dummy symbol $\perp$ otherwise.}.
%check the (\log 1/\beta) term
By integrating it into our framework, for all $t\in \{1,2,\dots,T\}$ simultaneously, we have:
\begin{itemize}
    \item With high probability $1-\beta$, the total error is bounded by: 
    \[
    O\Big(
    \frac{\log ^{1.5} T \cdot \log M \cdot m_{\max}(D^t)}{\varepsilon} \cdot (\log \frac{T \cdot \log M}{\beta}+\sqrt{\log\frac{1}{\delta}})
    \Big)
    \]
    \item The expected number of messages is $O( t\cdot M\log T)$.
\end{itemize}


% please check, there is already a paper talks about PDP in shuffle model


\section{Extension to Personalized DP}


% One interesting .. is PDP ... in this model ..., which alllows for 

% Here, we demonstrate that we can extend... to PDP, while still, roughly speaking achieves ... xxx


% What is epsilon min D? Why we want to achive it

% The first challenge is to formalise shuflle PDP with xxx

% definition 

% The second challenge is to dertermine epsilon min D

% We propose that this can be solved via a similar .... to that of m max


Our proposed framework is designed with high versatility, allowing it to be adapted to a broad range of DP settings beyond the standard definition. In this section, we extend our protocol to the Personalized Differential Privacy (PDP) setting to demonstrate its flexibility.



\paragraph{Problem Formulation}
The standard DP model provides uniform privacy protection for all users, which may not be desirable in real applications. In practice, different users may have heterogeneous privacy requirements, which is known as PDP in the literature~\cite{alaggan2015heterogeneous,ebadi2015differential,jorgensen2015conservative,sun2024personalized}.
In this setting, each user $i$ has an individual privacy budget $\varepsilon_i \in [\check{\varepsilon}, \hat{\varepsilon}]$, where $\hat{\varepsilon}$ and $\check{\varepsilon}$ denote the pre-defined global maximum and minimum privacy budgets across all users, respectively. 
For a dataset $D$, let \( \varepsilon_{\min}(D) = \min_{i \in [n]} \varepsilon_i \) denote the actual minimum privacy budget in the dataset. 


% here add a definition


\begin{definition}[Personalized Shuffle Differential Privacy] \label{def: PSDP} Let each user $i \in [n]$ have an individual privacy budget $\varepsilon_i \in [\check{\varepsilon}, \hat{\varepsilon}]$. A shuffle protocol $\mathcal{M}(D) = \mathcal{S} \circ \mathcal{R}(D)$ satisfies \emph{$(\{\varepsilon_i\}_{i=1}^n, \delta)$-personalized shuffle differential privacy} (PSDP) if for every user $i$, every fixed dataset $D_{-i}$ of the remaining users, and every subset of outputs $Y \subseteq \mathcal{Y}$: \[ \Pr[\mathcal{M}(D) \in Y] \leq e^{\varepsilon_i} \cdot \Pr[\mathcal{M}(\perp \uplus  D_{-i}) \in Y] + \delta, \] 
\end{definition}


Our goal is to design a mechanism that adapts any standard user-level shuffle-DP protocol to the PDP setting while achieving instance-specific utility that depends on $\varepsilon_{\min}(D)$ rather than the worst-case bound $\check{\varepsilon}$.




% please check: 这里的过渡用m_max(D)还是用central-DP下的做法？maybe the latter is better. or consider combination
It is clear that, the PDP problem shares a fundamental similarity with the user-level shuffle-DP problem addressed in previous sections. 
Recall that in the user-DP setting, users contribute different numbers of records $m_i \in [1, M]$, and our framework leverages domain partition to identify a threshold $m_\tau$ that approximates $m_{\max}(D)$, achieving instance-specific error. 
In the PDP setting, the heterogeneity shifts from record counts to privacy budgets. 
We apply an analogous domain partition strategy: just as we partition the contribution domain $[1, M]$ to find $m_\tau$, we now partition the privacy budget domain $[\check{\varepsilon}, \hat{\varepsilon}]$ to identify an instance-specific threshold $\varepsilon_\tau$ that approximates $\varepsilon_{\min}(D)$.
Consequently, the user-level PDP problem can be reduced to a standard user-level DP problem with a uniform privacy budget $\varepsilon_\tau$, allowing us to invoke our established framework as a modular building block.

Actually, this idea is already proposed by~\cite{chen2025general} for PrDP in central-DP at the record level. Our solution extends it to the user-level shuffle-DP setting. 

% \begin{algorithm}[t]
% \caption{Randomizer for Estimating $\varepsilon_\tau$}
% \label{alg: randomizer for estimate eps}
% \KwParam{$\varepsilon,\delta,n,M,\mathcal{P}_{Q_\text{count}}$\cite{ghazi2021differentially}$=(\mathcal{R}_\mathrm{count},\mathcal{S}_\mathrm{count},\mathcal{A}_\mathrm{count})$}
% \KwIn{$\varepsilon_i$}

% \For{$j \leftarrow 1$ \KwTo $\lceil \log(\hat{\varepsilon}/\check{\varepsilon})\rceil $ }
% {
% $b \leftarrow \mathbb{I}(\varepsilon_i \in (2^{j-1}\cdot \check{\varepsilon}, \min(2^j\cdot \check{\varepsilon},\hat{\varepsilon})])$;\tcp*{Estimate $\varepsilon_\tau$}

% $C_i^{(j)} \leftarrow \mathcal{R}_\mathrm{count}(b; \frac{2^{j-1}\cdot\check{\varepsilon}}{2}, \frac{\delta}{2}, n)$;

% \textbf{Send} $C_i^{(j)}$;
% }

% % \textbf{Run} Algorithm~\ref{alg:Randomizer for Estimating tau} $(m_i,\frac{\varepsilon}{2},\frac{\delta}{2},n,\mathcal{P}_\mathrm{cnt})$; \tcp*{Estimate $m_\tau$}

% \end{algorithm}

% \begin{algorithm}[t]
% \caption{Analyzer for Estimating $\varepsilon_\tau$}
% \label{alg: analyzer for estimate eps}
    
% $\varepsilon_\tau \leftarrow 0$;

% \For{$j \leftarrow \lceil \log(\hat{\varepsilon}/\check{\varepsilon})\rceil$ \KwTo $ 1$}{

% $\tilde{Q}_\mathrm{count}^{(j)} \leftarrow \mathcal{A}_\mathrm{count}(Z^{(j)}; \frac{2^{j-1}\cdot\check{\varepsilon}}{2}, \frac{\delta}{2},\frac{\beta}{2( \lceil \log(\hat{\varepsilon}/\check{\varepsilon})\rceil)},n)$;

% \If{$\tilde{Q}^{(j)}_\mathrm{count} > \frac{2}{2^{j-1}\cdot\check{\varepsilon}}\ln\frac{2\lceil \log(\hat{\varepsilon}/\check{\varepsilon})\rceil}{\beta}$}{
% $\varepsilon_\tau \leftarrow 2^{j-1}\cdot \check{\varepsilon}$;
% }

% }

% \Return $\varepsilon_\tau$;
% \end{algorithm}


\subsection{Static Setting}
We first discuss the static setting. 

\subsubsection{Multiple-Round Protocol}

Our solution consists of two phases:

(1) Identifying $\varepsilon_\tau$. 
We partition the privacy budget range $[\check{\varepsilon}, \hat{\varepsilon}]$ into $\lceil \log(\hat{\varepsilon}/\check{\varepsilon}) \rceil$ geometrically-spaced intervals: $(2^{j-1} \cdot \check{\varepsilon}, 2^j \cdot \check{\varepsilon}]$ for $j = 1, 2, \ldots, \lceil \log(\hat{\varepsilon}/\check{\varepsilon}) \rceil$. 
Each user determines which interval contains their privacy budget $\varepsilon_i$ and encodes this information as a binary indicator. 
We then apply a counting protocol from~\cite{ghazi2021differentially} with privacy budget scaled to the interval's lower bound to estimate the number of users in each interval. 
The analyzer identifies $\varepsilon_\tau$ as the smallest privacy budget threshold such that sufficiently many users have budgets above it (Algorithm~\ref{alg: randomizer for estimate eps} and~\ref{alg: analyzer for estimate eps}).

(2) Running the standard user-level shuffle-DP protocol. Once $\varepsilon_\tau$ is identified, we directly invoke our user-level shuffle-DP framework from previous sections with privacy budget $\varepsilon_\tau$.


\begin{theorem}
    % please check the 4/eps_min, you can do it in proof
    Given any $\varepsilon > 0, \delta >0, n \in \mathbb{Z}_+, $ $U \in \mathbb{Z}_{+}$, $M \in \mathbb{Z}_{+}$, for any query $Q$ and any base shuffle-DP protocol $\mathcal{P}_Q$, by integrating our two-round user-level shuffle-DP protocol to this framework, we have:
    \begin{itemize}
        \item The total error is bounded by:
        \[
            \begin{aligned}
            &O\Big(\gamma_{\ell_p}\big(Q,\frac{ m_{\max}(D)}{\varepsilon_{\min}(D)}\cdot \log \frac{\lceil\log(\hat{\varepsilon}/\check{\varepsilon})\rceil\cdot\log M}{\beta}  \big)\Big)
            +   \\
            &
            \err_{\ell_p}\Big(\mathcal{P}_Q, \frac{\varepsilon_{\min}(D)}{m_{\max}(D)},-, n \cdot m_{\max}(D),\frac{\log M}{\beta}\Big)
            \end{aligned}
        \]
        % please check this problem: when clip records corresponding to eps_tau, the max_(D) changes. so the bias and noise are not balance
        \item The expected number of message is $O(1)+\text{Msg}(\mathcal{P}_Q,\frac{\varepsilon_{\min}(D)}{m_{\max}(D)},-,n\cdot m_{\max}(D))$
    \end{itemize}

By integrating our one-round user-level shuffle-DP protocol to this framework, we have:

\begin{itemize}
    \item The total error is bounded by:

    % please check this problem: when clip records corresponding to eps_tau, the max_(D) changes. so the bias and noise are not balance
    \item The expected number of message is 
\end{itemize}
\end{theorem}


\subsubsection{Single-Round Protocol}
For all possible $\varepsilon_\tau$, integrate our one-round protocol.



% please check, here is a problem: you can not directly execute count query, because each user holds different m_i. 
% 之前的文章能直接clip，是因为一个人只有一条数据，总体被discard 的records是能被bound的。但在user-PDP setting下，如果很小的eps的人有很多的records呢？
% clip的原则是否是bound bias？eps很小且record很多的user值得关注吗？
% 注意到，user最多也就是m_i=m_max(D), 最终就是O(m_max(D)* 1/eps_min), 应该可以接受？noise里应该也是这个term。可以推一下



% 第二个问题，第一轮clip掉的是user...那调用user-level 算法的时候，n是unknown的。
% 解决：可以让被clip的用户正常参加，用dummy就行

% for pure DP, we easily ...., for (eps,delta)-dp, we can assume that delta is scale with epsilon.

% partition的范围从[1,M]到[xx,xx]

% 所有的error guarantee 都是把某个term换成某个term

% theorem 


% for the dynamic setting, similar to static setting, 其实非常一致。

% 只能用two-round实现，因为estimate m_\tau的时候要用统一的eps，所以必须先得到eps_\tau，才可以用统一的eps_tau进行estimate m_\tau和query evaluation

% delta? assume uniform delta or scale with eps? maybe the latter make sense






\subsection{Dynamic Setting}


% 现有数据还是先有人？eps是一开始就确定的嘛？dynamic setting下，eps会变吗？
It is clear that, similar to our user-level shuffle-DP protocol (Section~\ref{sec:Framework for Dynamic User-level Shuffle-DP}), we can adopt binary mechanism to address dynamic $\varepsilon_\tau$ estimation.




% 对于rank error, add one / change one 
\section{Experiment}

% if do median: consider not to use ziyue's protocol. you can make a short explanation that median can be reduce to range counting by xxx, so we directly use rm2..
\input{sec/9-Conclusion}


\clearpage



\bibliographystyle{ACM-Reference-Format}
\bibliography{ref}

\clearpage
\appendix
\renewcommand{\thetheorem}{4.\arabic{theorem}}
\setcounter{theorem}{0} 
\section{Proof of Theorem~\ref{the: static two-round}}

\label{sec:Proof of static two}


\begin{theorem} 
    Given any $\varepsilon > 0, \delta >0, n \in \mathbb{Z}_+, $ $U \in \mathbb{Z}_{+}$, $M \in \mathbb{Z}_{+}$, for any query $Q$ and any base shuffle-DP protocol $\mathcal{P}_Q$, our two-round protocol satisfies the following properties in the static setting:

    \begin{itemize}
    \item The messages received by the analyzer preserve $(\varepsilon,\delta)$-DP;
    \item The total error of this protocol is bounded by:
    \[
    \begin{aligned}
    &\mathrm{\gamma}_{\ell_p}\!\Big(
        Q,\;
        \frac{16}{\varepsilon}\ln\frac{2(\log M +1)}{\beta} \cdot m_{\max}(D)
    \Big)
    \\
    &+\;
    \err_{\ell_p}\!\Big(
        \mathcal{P}_Q,\;
        \frac{\varepsilon}{4m_{\max}(D)},\;
        \frac{\delta}{4m_{\max}(D)},\;
        \frac{\beta}{2},\;
        n\cdot 2m_{\max}(D)
    \Big)
    \end{aligned}
    \]



    \item In expectation, each user sends $1+ o(1)+\mathrm{Msg}(\mathcal{P}_Q, \allowbreak \frac{\varepsilon}{4m_{\max}(D)}, \allowbreak \frac{\delta}{4m_{\max}(D)}, \allowbreak n\cdot 2m_{\max}(D))$ messages.

\end{itemize}

\end{theorem}

\begin{proof}
    For privacy, the privacy budget $(\varepsilon, \delta)$ is split evenly between the $m_\tau$ estimation task and the query evaluation task. 
For the first round, invoking Lemma~\ref{lem: Bit Counting Protocol}, we have that for every $j \in \{0,1,2,\dots,\log M\}$, subdomain $[2^j+1,2^j]$ preserves $(\varepsilon/2,\delta/2)$-DP. Given that each $m_i$ only has an impact on a single subdomain $[2^j+1,2^j]$, 
parallel composition (Lemma~\ref{lem:Parallel Composition}) applies across the disjoint subdomains $\{[2^j+1,2^j]\}_{j\in\{0,1,2,\dots,\log M\}}$, ensuring this round satisfies $(\varepsilon/2,\delta/2)$-DP.
% it follows that the collection $\{[2^j+1,2^j]\}_{j\in\{0,1,2,\dots,\log M\}}$ preserves $(\varepsilon/2,\delta/2)$-DP by parallel composition (Lemma~\ref{lem:Parallel Composition}). 
For the second round, the protocol $\mathcal{P}_Q$ preserves $(\varepsilon/2,\delta/2)$-DP. Finally, by sequential composition (Lemma~\ref{lem:Sequential Composition}), our two-round protocol preserves $(\varepsilon,\delta)$-DP.
% please check that use lemma x.x or xx theory (lemma x.x)

For utility, let $Q(D)$ be the true query result on original dataset $D$, $Q(D^\tau)$ be the true query result on the standardized dataset $D^\tau$, and $\tilde{Q}(D^\tau)$ be our final noisy output. By the triangle inequality, the overall error $|\tilde{Q}(D^\tau) - Q(D)|$ can be decomposed into two parts:
$$|\tilde{Q}(D^\tau) - Q(D)| \le \underbrace{|Q(D^\tau) - Q(D)|}_{\text{Bias}} + \underbrace{|\tilde{Q}(D^\tau) - Q(D^\tau)|}_{\text{Noise}}.$$
\textit{Bounding the Bias}: 
% The bias is caused by users whose contribution $m_i$ exceeds $m_\tau$. 
The bias stems from the truncation of real records during the contribution standardization process. Let $m'$ denote the total number of records discarded from users with $m_i > m_\tau$. Since the bias $|Q(D^\tau) - Q(D)|$ is bounded by $\gamma_{\ell_p}( Q, m')$, our primary objective is to derive an upper bound for $m'$.
Let $C_j$ and $\tilde{C}_j$ denote the true and estimated number of users in the subdomain $I_j = [2^{j-1}+1, 2^j]$, respectively. Let $m_\tau = 2^l$. The algorithm selects $l$ as the largest index passing the threshold $T= \frac{2}{\varepsilon}\ln\frac{2(\log M +1)}{\beta}$. This implies that for all $j > l$, the estimated count $\tilde{C}_j$ did {not} pass the threshold, i.e., $\tilde{C}_j \le T$. 
Lemma~\ref{lem: Bit Counting Protocol} guarantees that for any such subdomain $I_j$ ($j > l$), with probability at least $1-\frac{\beta}{2(\log M +1)}$, the estimation error is bounded by $|\tilde{C}_j - C_j| \le T$. Combining these inequalities ($C_j \le \tilde{C}_j + |\tilde{C}_j - C_j|$), for any $j > l$, the true count is bounded by:
% Lemma~\ref{lem: Bit Counting Protocol} ensures that for any unselected subdomain $I_j$ ($j > l$), with probability at least $1-\frac{\beta}{2(\log M +1)}$, we have 
% \[
% |\tilde{C}_j-C_j| \le \frac{2}{\varepsilon}\ln\Big(\frac{2(\log M +1)}{\beta}\Big).
% \]
% Consequently, for any $j > l$, the true count is bounded by:
$$ C_j \le \frac{4}{\varepsilon}\ln\Big(\frac{2(\log M+1)}{\beta}\Big).
% = O\Big(\frac{1}{\varepsilon}\ln\frac{\log M}{\beta}\Big).
$$
Now, we calculate the total number of excluded records $m'$. A user in subdomain $I_j$ (where $j > l$) contributes at most $2^j$ records but is clipped to $m_\tau = 2^l$. The loss is at most $2^j - 2^l$.
Note that for $j > K = \lceil \log m_{\max}(D) \rceil$, the true count $C_j = 0$. Thus, summing over $j$ from $l+1$ to $K$:
$$\begin{aligned}
m'
&= \sum_{j=l+1}^{K} C_j \cdot (2^j - 2^l) \\
&\le \frac{4}{\varepsilon}\ln\frac{2(\log M +1)}{\beta} \cdot \sum_{j=l+1}^{K} (2^j - 2^l) \quad  \\
% &\le \frac{4}{\varepsilon}\ln(\frac{\log M+2}{\beta}) \cdot \sum_{j=l+1}^{K} 2^j \\
% &= \frac{4}{\varepsilon}\ln(\frac{\log M+2}{\beta}) \cdot (2^{K+1} - 2^{l+1})\\
% &\le \frac{4}{\varepsilon}\ln\frac{2(\log M +1)}{\beta} \cdot 4 m_{\max}(D)\\
% \\
% & = O\left(m_{\max}(D)\cdot\frac{1}{\varepsilon}\log\frac{\log M}{\beta}\right) 
&\le \frac{4}{\varepsilon}\ln\frac{2(\log M +1)}{\beta} \cdot 2^{K+1} \\
&\le \frac{16}{\varepsilon}\ln\frac{2(\log M +1)}{\beta} \cdot m_{\max}(D).
\end{aligned}$$
This explicitly validates the property (a) claimed in section~\ref{sec: static round 1}, confirming that the number of excluded records is bounded by $\tilde{O}(m_{\max}(D)/\varepsilon)$.
Consequently, the bias is bounded by: $$\mathrm{\gamma}_{\ell_p}\big(Q, \frac{16}{\varepsilon}\ln\frac{2(\log M +1)}{\beta} \cdot m_{\max}(D)\big).$$
with probability $1-\frac{\beta}{2}$ by union bound.



\textit{Bounding the Noise:} 
The noise magnitude depends on the threshold $m_\tau$. Specifically, the base protocol $\mathcal{P}_Q$ is invoked with privacy budget $(\frac{\varepsilon}{2m_\tau},\frac{\delta}{2m_\tau})$ and dataset size $n \cdot m_\tau$. To bound the noise, it suffices to prove that with high probability, $m_\tau \le 2m_{\max}(D)$.
Still using $K = \lceil \log m_{\max}(D) \rceil$, for any subdomain index $j > K$, the subdomain $[2^{j-1}+1, 2^j]$ strictly exceeds $m_{\max}(D)$, implying the true count $C_j=0$. By Lemma~\ref{lem: Bit Counting Protocol}, the estimated count $\tilde{C}_j$ is pure noise. Therefore, with probability at least $1-\frac{\beta}{2(\log M +1)}$, $\tilde{C}_j$ does not exceed the error bound $T$.
% please check if it exactly 1-beta/2
By union bound, with probability at least $1-\beta/2$, the algorithm will not select any $j > K$. Therefore, the selected threshold satisfies $m_\tau \le 2^K \le 2 m_{\max}(D)$.
{This confirms property (b).}
% By lemma~\ref{lem: Bit Counting Protocol}, if a subdomain $[2^{j-1}+1,2^j]$ is empty, i.e., $C_j=0$, with probability $1-\frac{\beta}{2(\log M +1)} $, the estimated count $\tilde{C}_j \le \frac{2}{\varepsilon}\ln \frac{2(\log M +1)}{\beta}$. Thus, by union bound, with probability at least $1-\beta/2$, for all $j> m_\tau$, $[2^{j-1}+1,2^j] $ is empty. Therefore, we can guarantee that with probability at least $1-\beta/2$, $m_\tau\le 2\cdot m_{\max}(D)$. As a result, the noise caused by $\mathcal{P}_Q$ is bounded by $\err_{\ell_p}(\mathcal{P}_Q,\frac{\varepsilon}{4m_{\max}(D)},\frac{\delta}{4m_{\max}(D)},\frac{\beta}{2},   n\cdot 2m_{\max}(D))$ with probability at least $1-\beta/2$. 
Substituting this bound into the error function, the noise is bounded by $\err_{\ell_p}(\mathcal{P}_Q,\frac{\varepsilon}{4m_{\max}(D)},\frac{\delta}{4m_{\max}(D)},\frac{\beta}{2}, n\cdot 2m_{\max}(D))$ with probability at least $1-\frac{\beta}{2}$.

Combining bias and noise yields the final error bound:
    \[
    \begin{aligned}
    &\mathrm{\gamma}_{\ell_p}\!\Big(
        Q,\;
        \frac{16}{\varepsilon}\ln\frac{2(\log M +1)}{\beta} \cdot m_{\max}(D)
    \Big)
    \\
    &+\;
    \err_{\ell_p}\!\Big(
        \mathcal{P}_Q,\;
        \frac{\varepsilon}{4m_{\max}(D)},\;
        \frac{\delta}{4m_{\max}(D)},\;
        \frac{\beta}{2},\;
        n\cdot 2m_{\max}(D)
    \Big)
    \end{aligned}
    \]
    with probability at least $1-\beta$.

    For the communication cost, recall that our protocol consists of two rounds, each consuming privacy parameters $(\varepsilon/2,\delta/2,\beta/2)$. 
    In the first round, each user executes the bit-counting protocol \(\mathcal{P}_{Q_\text{count}}\)~\cite{ghazi2021differentially} on each subdomain $I_j$ and parallel composition (Lemma~\ref{lem:Parallel Composition}) applies across the \(\log M+1\) disjoint subdomains.
    Crucially, since these subdomains are disjoint, each user contributes a value of $1$ to exactly one subdomain, which we denoted as $I^*$,  and $0$ to all remaining $\log M$ subdomains.
    According to Lemma~\ref{lem: Bit Counting Protocol}, the expected message cost for $I^*$ is $1 + O(\frac{\log(1/\delta)}{\varepsilon n})$, while when input value is $0$, the protocol only sends out $o(\frac{\log(1/\delta)}{\varepsilon n})$ message.
    Consequently, the total expected communication per user in the first round is:
    \[
        1 \cdot \Big(1 + O\big(\frac{\log(1/\delta)}{\varepsilon n}\big)\Big) + \log M \cdot O\big(\frac{\log(1/\delta)}{\varepsilon n}\big) = 1 + o(1),
    \]
    assuming $n \gg \log M \cdot \log(1/\delta)/\varepsilon$.
    % please check if this inequation is right
    
    % By lemma~\ref{lem: Bit Counting Protocol}, for each subdomain, the expected number of messages sent by a user is bounded by $1+O(\frac{\log (1/\delta)}{\varepsilon n}) = O(1)$. Since there are \(O(\log M)\) subdomains, the total communication per user in this round is $O(\log M)$. 
    In the second round, the base protocol \(\mathcal{P}_Q\) is executed on the standardized dataset with privacy budget $(\frac{\varepsilon}{2m_\tau},\frac{\delta}{2m_\tau})$. Thus, the communication cost is then $\mathrm{Msg}(\mathcal{P}_Q,\frac{\varepsilon}{2m_\tau},\frac{\delta}{2m_\tau},n\cdot m_\tau)$. From the utility analysis, we have \(m_\tau \le 2 m_{\max}(D)\), 
    % please check, with probability xx?
    and therefore this cost is bounded by $\mathrm{Msg}(\mathcal{P}_Q,\frac{\varepsilon}{4 m_{\max}(D)},\frac{\delta}{4 m_{\max}(D)},n \cdot 2m_{\max}(D))$. Combining these two rounds, the total number of messages per user is bounded by:
    $$1+o(1)+ \mathrm{Msg}(\mathcal{P}_Q, \allowbreak \frac{\varepsilon}{4m_{\max}(D)}, \allowbreak \frac{\delta}{4m_{\max}(D)}, n\cdot 2m_{\max}(D)).$$


\end{proof}

\section{Proof of Theorem~\ref{the: static one-round}}

\setcounter{theorem}{1}
\begin{theorem}
    Given any $\varepsilon > 0, \delta >0, n \in \mathbb{Z}_+, $ $U \in \mathbb{Z}_{+}$, $M \in \mathbb{Z}_{+}$, under the static setting, our one-round protocol achieves that:

    \begin{itemize}
    \item The messages received by the analyzer preserve $(\varepsilon,\delta)$-DP;
    \item The total error of this protocol is bounded by:
    \[
    \begin{aligned}
    &\mathrm{\gamma}_{\ell_p}\!\Big(
        Q,\;
        \frac{16}{\varepsilon}\ln\frac{2(\log M+1)}{\beta} \cdot m_{\max}(D)
    \Big)
    +\;
    O\Big(\err_{\ell_p}\!(
        \mathcal{P}_Q,\;    \\
    &
        \frac{\varepsilon}{m_{\max}(D)\cdot \log M},\;
        \frac{\delta}{m_{\max}(D)\cdot \log M},\;
        \frac{\beta}{\log M},\;
        n\cdot m_{\max}(D)
    )\Big)
    \end{aligned}
    \]



    \item In expectation, each user sends $1+o(1)+\sum_{j \in \{0,1,\dots, \log M\}} \allowbreak \mathrm{Msg}(\mathcal{P}_Q, \frac{\varepsilon}{2\cdot 2^j(\log M+1)}, \frac{\delta}{2\cdot 2^j(\log M+1)},n \cdot 2^j)$ messages.

\end{itemize}

\end{theorem}

\begin{proof}
     For privacy, the privacy budget $(\varepsilon, \delta)$ is split evenly between the $m_\tau$ estimation task and the query evaluation task.
    For the $m_\tau$ estimation,since the subdomains $\{I_j\}$ are disjoint, parallel composition (Lemma~\ref{lem:Parallel Composition}) guarantees that this step satisfies $(\varepsilon/2, \delta/2)$-DP.
    For query evaluation, by applying sequential composition (Lemma~\ref{lem:Sequential Composition}) over the $\lceil \log M \rceil + 1$ candidate query results, each with budget $(\frac{\varepsilon}{2(\log M+1)}, \frac{\delta}{2(\log M+1)})$, the total cost sums to $(\varepsilon/2, \delta/2)$-DP.
    By basic composition of these two parts, the overall protocol satisfies $(\varepsilon, \delta)$-DP.

    For utility, similar to the proof of Theorem~\ref{the: static two-round}, the total error is bounded by the sum of the clipping bias and the DP noise.
    Observe that the procedure for determining $m_\tau$ (including the privacy budget allocation and the threshold condition) is algorithmically equivalent to the first round of the two-round protocol. Consequently, the analysis for the number of excluded records holds directly. With probability at least $1-\beta/2$, the clipping bias is bounded by:
    \[
    \mathrm{\gamma}_{\ell_p}\!\Big(
        Q,\;
        \frac{4}{\varepsilon}\ln(\frac{2(\log M+1)}{\beta}) \cdot 4 m_{\max}(D)
    \Big)
    \]
    Regarding the selected threshold, the guarantee from Theorem~\ref{the: static two-round} remains valid: with probability at least $1-\beta/2$, we have $m_\tau \le 2m_{\max}(D)$.
    However, unlike the two-round protocol, the privacy budget for query evaluation is split among $\log M + 1$ candidate subdomains. For the selected query result corresponding to $m_\tau$, the base protocol $\mathcal{P}_Q$ is applied with parameters:
    \[
    \varepsilon' = \frac{\varepsilon}{2(\log M+1) m_\tau}, \quad \delta' = \frac{\delta}{2(\log M+1) m_\tau}, \quad n' = n \cdot m_\tau.
    \]
    Substituting the upper bound $m_\tau \le 2m_{\max}(D)$ into the error expression $\mathrm{Error}_{\ell_p}(\mathcal{P}_Q, \varepsilon', \delta', \beta', n')$, the noise is bounded by:
    \[
    \begin{aligned}
        O\Big( \mathrm{Error}_{\ell_p}\!\big(
        &\mathcal{P}_Q,\;
        \frac{\varepsilon} {m_{\max}(D) \cdot \log M},\;
        \frac{\delta}{m_{\max}(D) \cdot\log M},\; \\
        &\frac{\beta}{\log M},\;
        n \cdot m_{\max}(D)
    \big) \Big).
    \end{aligned}
    \]
    Combining the bias and noise bounds yields the final overall error bound.

For the communication cost, our one-round protocol executes two tasks in parallel: $m_\tau$ estimation and query evaluation across all candidate thresholds.

For the $m_\tau$ estimation task, the analysis is identical to the first round of the two-round protocol. Each user executes the bit-counting protocol $\mathcal{P}_{Q_\text{count}}$~\cite{ghazi2021differentially} on each subdomain $I_j$ with budget $(\varepsilon/2, \delta/2)$. Since the subdomains are disjoint, each user contributes a value of $1$ to exactly one subdomain and $0$ to all others. By Lemma~\ref{lem: Bit Counting Protocol}, the expected message cost is:
\[
1 \cdot \Big(1 + O\big(\frac{\log(1/\delta)}{\varepsilon n}\big)\Big) + \log M \cdot O\big(\frac{\log(1/\delta)}{\varepsilon n}\big) = 1 + o(1),
\]
assuming $n \gg \log M \cdot \log(1/\delta)/\varepsilon$.

For the query evaluation task, each user must send messages for all $\log M + 1$ candidate thresholds $\{2^j\}_{j=0}^{\log M}$. For each subdomain $j$, the user constructs a standardized dataset $D’$ of size $2^j$ and applies the base randomizer $\mathcal{R}_Q$ with privacy budget $(\frac{\varepsilon}{2 \cdot 2^j (\log M+1)}, \frac{\delta}{2 \cdot 2^j (\log M+1)})$ and participation size $n \cdot 2^j$. The communication cost for subdomain $j$ is therefore $\mathrm{Msg}(\mathcal{P}_Q, \frac{\varepsilon}{2\cdot 2^j(\log M+1)}, \frac{\delta}{2\cdot 2^j(\log M+1)}, n \cdot 2^j)$.

Summing over all $\log M + 1$ subdomains, the total expected communication per user is:
\[
1 + o(1) + \sum_{j=0}^{\log M} \mathrm{Msg}\Big(\mathcal{P}_Q, \frac{\varepsilon}{2\cdot 2^j(\log M+1)}, \frac{\delta}{2\cdot 2^j(\log M+1)}, n \cdot 2^j\Big).
\]

\end{proof}

   
\section{Proof of Theorem~\ref{the: dynamic setting guarantee}}

\renewcommand{\thetheorem}{5.\arabic{theorem}}
\setcounter{theorem}{0} 
\begin{theorem}
    Given any $\varepsilon > 0, \delta > 0$, any $T \in \mathbb{Z}_+$, and $n \in \mathbb{Z}_+$, $M \in \mathbb{Z}_+$, for any query $Q$, at all timestamps $t \in \{1, 2, \dots, T\}$ simultaneously, our protocol for the dynamic setting achieves the following guarantees:
\begin{itemize}
    \item The messages received by the analyzer preserve $(\varepsilon,\delta)$-DP;
    \item With probability at least $1-\beta$, the total error is bounded by:
    \[
    \begin{aligned}
    &O\Big(\gamma_{\ell_p}\big( Q, \frac{\log ^{1.5} T \cdot \log M}{\varepsilon} \ln \frac{T \cdot \log M}{\beta} \cdot m_{\max}(D^t)\big)\Big) + O\Big(\err_{\ell_p}\big(\mathcal{P}_Q,
    \\
    & \frac{\varepsilon}{m_{\max}(D^t) \cdot \log M}, \frac{\delta}{m_{\max}(D^t) \cdot \log M}, \frac{\beta}{\log M}, n \cdot m_{\max}(D^t)\big)\Big)
    \end{aligned}
    \]
    \item In expectation, each user sends $(\log M+1)(1+o(1)) + \sum_{j=0}^{\log M} \sum_{k=1}^{2^j} \mathrm{Msg}(\mathcal{P}_Q, \frac{\varepsilon}{2\cdot 2^j(\log M+1)}, \frac{\delta}{2\cdot 2^j(\log M+1)}, n \cdot 2^j)$ messages over the entire time horizon $T$.
\end{itemize}
\end{theorem}


\begin{proof}
    For privacy, the privacy budget $(\varepsilon, \delta)$ is split evenly between the two parallel tasks: dynamic counting and query evaluation, each receiving $(\varepsilon/2, \delta/2)$.

For \emph{dynamic counting}, the user participates in $\log M+1$ interval counting mechanisms. Since the intervals $\{I_j\}$ are disjoint (a user belongs to exactly one interval at any time), we apply parallel composition (Lemma~\ref{lem:Parallel Composition}) over the intervals, allocating budget $(\varepsilon', \delta') = (\frac{\varepsilon}{2(\log M+1)}, \frac{\delta}{2(\log M+1)})$ per interval.
Within each interval $j$, the binary mechanism operates over $\log T+1$ tree levels. By applying sequential composition (Lemma~\ref{lem:Sequential Composition}) over these levels, each level receives budget $(\frac{\varepsilon'}{\log T+1}, \frac{\delta'}{\log T+1})$.
Thus, dynamic counting satisfies $(\varepsilon/2, \delta/2)$-DP.

For \emph{query evaluation}, each user processes their data under all $\log M+1$ candidate thresholds. These thresholds are not disjoint (threshold $j+1$ covers threshold $j$), requiring sequential composition over the $\log M+1$ intervals with budget $(\varepsilon', \delta')$ per interval.
For each interval $j$, the user contributes $2^j$ streams, each running the base protocol with budget $(\frac{\varepsilon'}{2^j}, \frac{\delta'}{2^j})$. By composition over the $2^j$ streams, interval $j$ satisfies $(\varepsilon', \delta')$-DP.
Under the continual observation model~\cite{dwork2010continual,chan2011private}, incremental updates to streams incur no additional privacy cost, as the guarantee applies to the user's complete sequence $D_i = \{x_{i,t}\}_{t=1}^T$.
Thus, query evaluation satisfies $(\varepsilon/2, \delta/2)$-DP.

By basic composition of the two tasks, the overall protocol satisfies $(\varepsilon, \delta)$-DP.


For utility, similar to the static one-round protocol, the total error is bounded by the sum of clipping bias and DP noise.



\emph{Clipping Bias.}
Similar to the proof of Theorem~\ref{the: static two-round}, by Lemma ..., at each time $t$, the true count of interval $I_j$ is bounded by:
\[
C_j \le O(\frac{\log^{1.5} T}{\varepsilon'}\ln \frac{T}{\beta'}).
\]
Let $l = j^*$ and $K = \lceil \log m_{\max}(D^t) \rceil$. We can get the total number of discarded records is:
\[
\begin{aligned}
m'
&= \sum_{j=l+1}^{K} C_j \cdot (2^j - 2^l) \\
&\le O\Big(\frac{\log ^{1.5} T}{\varepsilon'} \ln \frac{T}{\beta'}\Big) \cdot \sum_{j=l+1}^{K} (2^j - 2^l) \\
&\le O\Big(\frac{\log ^{1.5} T \cdot \log M}{\varepsilon} \ln \frac{T\log M}{\beta}\Big) \cdot m_{\max}(D^t).
\end{aligned}
\]




\emph{DP Noise.}
For the selected threshold $m_\tau^t = 2^{j^*}$, the base protocol $\mathcal{P}_Q$ is applied with parameters:
\[
\varepsilon' = \frac{\varepsilon'}{2^{j^*}} = \frac{\varepsilon}{2(\log M+1) \cdot 2^{j^*}}, \quad \delta' = \frac{\delta}{2(\log M+1) \cdot 2^{j^*}}, \quad n' = n \cdot 2^{j^*}.
\]
Substituting the upper bound $2^{j^*} \le 2m_{\max}(D^t)$, the noise is bounded by:
\[
\begin{aligned}
    O\Big( \err_{\ell_p}\big( &\mathcal{P}_Q, \frac{\varepsilon}{m_{\max}(D^t) \cdot \log M}, \frac{\delta}{m_{\max}(D^t) \cdot \log M}, \\ &\frac{\beta}{\log M}, n \cdot m_{\max}(D^t)\big) \Big).
\end{aligned}
\]

Combining the bias and noise bounds yields the stated error guarantee.

For communication cost, the protocol executes two tasks in parallel: dynamic counting and query evaluation.

For \emph{dynamic counting}, each user runs $\log M+1$ independent binary mechanism instances (one per interval). 
Since intervals are disjoint, each user contributes to exactly one interval at any given time. 
By Lemma~\ref{lem: Bit Counting Protocol}, each binary mechanism instance with parameters $(\frac{\varepsilon'}{\log T+1}, \frac{\delta'}{\log T+1})$ sends $1 + o(1)$ messages in expectation.
Over $\log M+1$ intervals, the total cost is:
\[
(\log M+1) \cdot (1 + o(1)) = (\log M+1)(1+o(1)).
\]

For \emph{query evaluation}, the user maintains $\sum_{j=0}^{\log M} 2^j = 2^{\log M+1}-1$ streams across all intervals.
Each stream $S[j][k]$ runs an instance of the base protocol $\mathcal{R}_Q$ with budget $(\frac{\varepsilon'}{2^j}, \frac{\delta'}{2^j})$ and participant count $n \cdot 2^j$.
The communication cost for stream $S[j][k]$ is $\mathrm{Msg}(\mathcal{P}_Q, \frac{\varepsilon'}{2^j}, \frac{\delta'}{2^j}, n \cdot 2^j)$.
Summing over all streams:
\[
\sum_{j=0}^{\log M} \sum_{k=1}^{2^j} \mathrm{Msg}\Big(\mathcal{P}_Q, \frac{\varepsilon}{2\cdot 2^j(\log M+1)}, \frac{\delta}{2\cdot 2^j(\log M+1)}, n \cdot 2^j\Big).
\]

Combining both tasks, the total expected communication per user is:
\[
\begin{aligned}
    &(\log M+1)(1+o(1)) \\&+ \sum_{j=0}^{\log M} \sum_{k=1}^{2^j} \mathrm{Msg}\Big(\mathcal{P}_Q, \frac{\varepsilon}{2\cdot 2^j(\log M+1)}, \frac{\delta}{2\cdot 2^j(\log M+1)}, n \cdot 2^j\Big).
\end{aligned}
\]

\end{proof}
\end{document}